\chapter{Anhang}
\section{Interview}
\label{app:interview}
Transkript
26. August 2025, 12:00PM

Teich Christian (PS-TC/GP)   0:10
So hello, young Christian here. Hi, check, check. It's to meet you on. I'm sorry.
Unmute.

Yan Ohainski   0:21
Yeah, it's sorry. Hello, mute. Yeah, inside.

Teich Christian (PS-TC/GP)   0:25
Yeah.
Yeah, Ganne, uh, the roof is getting another roof one.

Yan Ohainski   0:31
You know, it's been getting look on Anika's so meeting.
And.

Teich Christian (PS-TC/GP)   1:01
Dinge oder Antworten auf gestellte Fragen liefert ne, weil das, was wir derzeit erleben, ist genau in der Richtung. Das ist sehr hemdsärmelig, das ist sehr, sag ich mal an den User und es ist auch revolutionär da braucht man nicht viel Fantasie, dass das auch umsetzbar ist, aber das wäre ein sehr.

Yan Ohainski   1:07
Ja.
Hm.

Teich Christian (PS-TC/GP)   1:21
Guter Move, um deutlich besser zu werden in der Kommunikation. Vielleicht gibst du mir aber trotzdem noch mal background, woher kommst du du hast irgendwie schreibst eine Masterarbeit habe ich verstanden gehabt und machst das dann innerhalb des Startups oder wie ist das?

Yan Ohainski   1:22
Mhm.
Ja.
Ja klar.
Genau ich wollte ihnen sowieso ganz kurz noch mal n Abriss geben. Was macht Noticent und dann mein Projekt noch mal erklären? In dem Zuge auch noch mal wer bin ich, was mach ich und wieso mach ich das ja, aber ich muss noch einmal oh, ich kann leider nichts teilen, ich hab so ne Präsentation vorbereitet wenn.

Teich Christian (PS-TC/GP)   1:42
Ja.
Mhm.
Ah, ja genau nee, weil ich eingeladen hab stimmt ich muss den immer freigeben kleinen Moment, das ist so ein bisschen irgendwie so ein Sicherheitsding gerade mal frei kleinen Moment, so Status in Referenz ändern.

Yan Ohainski   1:59
Ja, auch auch verständlich.
Solange es am Ende geht.

Teich Christian (PS-TC/GP)   2:07
So.

Yan Ohainski   2:09
Ja, und du hattest ja auch die Aufzeichnung schon gestartet, wie ich richtig sehe, aber sehr gut ja.

Teich Christian (PS-TC/GP)   2:09
So jetzt.
Ganz genau ne, ich wollte mal gleich anhand dessen also wir machen das gleich als live Erlebnis und da bin ich schon sage ich mal da mache ich mehr als meine Kollegen, die ich dazu bewegen will. Wir sind alle so um die in dem Alter wie du jetzt siehst junger Mann helfen sie mir mal über die Straße, ne, das ist typisch Management, die sind.

Yan Ohainski   2:18
Ja.
Ah, okay, ja.

Teich Christian (PS-TC/GP)   2:31
Alle älter und wir üben uns zumindest damit in verbesserter Dokumentation und der Grund dahinter als Kontext auch gegeben und als Motivation und dann steigen wir bei euch ein ist wir haben mittlerweile sehr schlanke Organisationen über viele Hierarchie.

Yan Ohainski   2:44
Mhm.

Teich Christian (PS-TC/GP)   2:50
Stufen hinweg ja, haben wir uns sehr gesund geschrumpft? Dadurch sind wir wettbewerbsfähig geworden, aber es werden viele Dinge kurz mit irgendwem besprochen und leider sehr schlecht dokumentiert. Ne und das führt dazu ich habe die Verantwortung für einen Bereich, der hat in.

Yan Ohainski   3:05
Ach ja.

Teich Christian (PS-TC/GP)   3:10
Umsatz 1,5 Milliarden Euro also, das ist schon ein größerer Bereich und ganz häufig erwische ich mich selber, aber auch meine Kollegen, dass ich denen dann sagen muss das hatte ich doch schon mit denen besprochen, das hatten wir doch schon entschieden, oder? Das haben wir schon gemacht und dann erkennt man sich.

Yan Ohainski   3:14
Ja ja.

Teich Christian (PS-TC/GP)   3:27
Ja kleinen Moment, ich guck mal in meinen Mails wo natürlich habe ich nie eine Mail rausgeschickt natürlich habe ich nie, irgendwann hat man das mal besprochen und schon ist man in dieser Dokumentationsfalle und ein Trick, um da besser zu werden. Ist die Zeit, dass wir in der Tat die Chance nutzen, hier über ein Transkript.

Yan Ohainski   3:32
Ja.

Teich Christian (PS-TC/GP)   3:45
Mitzuschreiben in MS Teams. Deswegen musst du es doppelt durchklicken es wird also alles aufgezeichnet und dann lass ich das hinterher. Über eine AI laufen hier in dem Meeting wäre es jetzt nicht wichtig, aber ich würde das mal exemplarisch dann ausprobieren, wie das so mit unseren Inhalten läuft.

Yan Ohainski   3:47
Mhm.
Ja.
Mhm.
Ja ausprobieren wollen, ne?
Hm.

Teich Christian (PS-TC/GP)   4:04
Dann gehe ich in der Regel noch mal drüber, weil ich sag mal so 20 bis 30% der Inhalte müssen geändert werden, sonst sind sie nicht korrekt oder eine nicht hundertprozentig präzise und dann werden die versendet und dann ist das Problem behoben. Allerdings dauert das alles noch sehr lange.

Yan Ohainski   4:13
Ja.
Hm.

Teich Christian (PS-TC/GP)   4:24
Ne und dann kommen wir gleich nachher zu den Paint Points und was man da vielleicht machen könnte.

Yan Ohainski   4:30
Ja ist klar, das wäre auch eine meiner Fragen gewesen wie sieht denn dein Workflow aus? Aber das hast du mir jetzt schon vorweggenommen schon mal gut nee und dann starte ich erst mal erzähl mal was über noti Cent selbst und dann kommen wir auch zu einem heißen Thema. Erst mal nooti Cent, selbst nooti Cent selbst ist ein sicherer KI Assistent, der die informations.

Teich Christian (PS-TC/GP)   4:35
Genau.

Yan Ohainski   4:47
Flut reduzieren soll und die Kommunikation halt transformiert, so wie es dort steht. Genau da werden mehrere Bereiche angebunden, an den Assistenten selbst die Kontakte, der Kalender und to do's Aufgaben genauso wie eine zentralisierte Suche, in der man die Möglichkeit hat, auf alle seine Daten, die ich gerade genannt hatte, die Kontakte, Kalenderaufgaben und Mails zuzugreifen.
Und der KI Assistent behält das Ganze im Überblick und kann dann so basierend auf Inputs, die der User gibt, bereits schon vorher selektieren. Was könnte relevant sein? Was ist wichtig, was ist nicht so wichtig für dieses Gespräch für diese Mail?
Um so halt einfach das lästige suchen und herauszufinden OK wo stand das noch mal zu vermeiden?

Teich Christian (PS-TC/GP)   5:28
By the way schickst du hinterher die Folien kann ich die haben oder ist das schwierig?

Yan Ohainski   5:29
Ja.
Kann an sich, hab ich da gar kein Problem mit ich frag noch mal die anderen die ja was zu sagen haben und dann kriegen wir das sicher hin, also das sind ja auch.

Teich Christian (PS-TC/GP)   5:43
Ich behandle, ich behandle die vertraulich. Aber würde sie in unserem Führungskreis ne, das ist so das Board of Management zeigen, weil ich gerne für die Zusammenarbeit in dem Zusammenhang werben möchte ne, ich würde das vermutlich nicht selber machen, da gebe ich dann schnell ab.

Yan Ohainski   5:46
Ja klar.
Ja.
Ja.
Ja.

Teich Christian (PS-TC/GP)   6:03
Aber dass dafür geworben werden kann, und da wäre es hilfreich, wenn ich das zeigen könnte.

Yan Ohainski   6:08
Ja, absolut verstehe ich auch voll, ich kann es leider selbst nicht entscheiden, einfach so deswegen würde ich noch mal nachhaken, aber ich denke, da spricht absolut nichts gegen das sind auch die Folien, die ich hier jetzt gerade gezeigt habe und die nächsten 23 die noch kommen, sind auch Folien von Luca und Anika selbst, die sie für genau sowas entwickelt haben, um sie halt potenzielle Investoren zu zeigen. Also das sind eigentlich genau welche dafür.

Teich Christian (PS-TC/GP)   6:26
Mhm.

Yan Ohainski   6:29
Ja, natürlich kann ich dir den zuschicken genau ich hatte über den KY stand gesprochen und falsche Taste so ist halt die Frage wie funktioniert das Ganze? Und zwar soll das System nötig sind an eure Systeme oder bestehende Systeme angebunden werden. Hier in dem Fall Microsoft 365 es gibt noch andere Gmail Motion.
Lessieren also, was man so in der Industrie findet und dadurch bekommt der Agent No decent die Möglichkeit, darauf zuzugreifen und sie aufzubereiten und hat dann halt die Möglichkeit, Aufgaben daraus zu ziehen, die Mail zu priorisieren, an sich auch wissen in Lexikas oder Glossaren, die es gibt. In der im Unternehmen halt darauf zuzugreifen und weitere.
Und aktuell ist das Ganze noch als Add in für Outlook geplant, hauptsächlich auf der nächsten Seite. So soll das Aussehen, das ist jetzt noti sellet Assistent allgemein noch nicht mit dem Meeting, das kommt gleich und hier ist dann halt an der Seite n extra Reiter, der Halt zusätzlich gebündelt Informationen anzeigt. Basierend auf der Mail, die hier bereits angezeigt wird.
Und welchen die im Zusammenhang damit stehen?

Teich Christian (PS-TC/GP)   7:31
Mhm.

Yan Ohainski   7:32
Genau, und hier rechts sieht es auch noch mal was konkret möglich ist, da würde ich jetzt aber hier nicht so im Detail drauf eingehen, weil es nicht Teil meiner Aufgabe ist, sozusagen.
Und dazu darum wird jetzt der Meeting Assistent als eigenes Produkt in das gesamt Constructor Environment noti Cent eingeführt und das Ganze sieht so aus im Workflow aber erstmal würde ich vorweg noch was zu mir sagen, bevor ich da jetzt weiter einsteige und weshalb ich das mache also ja ist richtig, ich schreib meine Masterarbeit.
Jetzt über das Thema Ich hatte vorher auch schon mit Notisens zusammen an Speech to Text, also Sprachumwandlung in Textform gearbeitet. Dort versucht in den Chat von Notisens selbst eine Spracheingabefunktion einzubauen und dort dann auch ein bisschen Aufnahmen gemacht und trainiert und die KI verbessert, sodass man dann auch wirklich sprechen kann.
Und darauf aufbauend war dann die Idee von Luca und Henry, die sich technisch um die Sachen kümmern Mensch, machen wir daraus doch meetingassistent, das ist etwas, was Mehrwert bietet, was es noch nicht in der Form so gibt vielleicht abgewandelt und da wir uns alle seit dem Studium schon kennen, also im Master gemeinsam alle angefangen hier.
Dachte sich OK ja, das wär doch super für dich. Wir ziehen wir haben was davon das Produkt am Ende und ich hab am Ende ne Abschlussarbeit mit einem Thema wo ich weiß, das wird auch verwendet. Das ist nicht nur einfach nur damit ich beschäftigt werde von meinem Prof und so kam das Ganze dann zustande und an sich studiere ich auch elektroninformationstechnik wie Luca.

Teich Christian (PS-TC/GP)   8:56
Ah ja, sehr schön Mhm.

Yan Ohainski   8:57
Genau also, ich komme eher aus dem technischen Bereich, weshalb ich auch sehr an dem Gespräch interessiert war, weil ich kenne das nur aus dem technischen, also für mich ist immer wichtig OK. Wir müssen jedes ganze Detail auffassen und es muss wirklich genau sein und auch sehr genau beschrieben sein, nichts verallgemeinert sondern.
Und Sorry Luca meinte da, dass du ja auch Punkte hast aus dem Managementbereich, die ja vielleicht nicht so interessiert am technischen sind.

Teich Christian (PS-TC/GP)   9:22
Mhm.

Yan Ohainski   9:25
Ja, und so soll das Ganze dann aussehen, also hier einmal ganz kurz stark vereinfacht in den User Links zu sehen, der spricht erstmal einfach in das System rein, da speech to Text hatte ich mich ja bereits drum gekümmert und unten in dem Lilanen siehst du dann meist siehst du dann hier ein Beispiel wie das ganze aussehen soll?

Teich Christian (PS-TC/GP)   9:26
Ah.

Yan Ohainski   9:41
Das Gesprochene wird vorher von einer weiteren KI analysiert, die dann quasi die ersten Stichpunkte rauszieht. Sieht man hier in dem Reiter Voranalyse aus dem Text OK. Dann schickt Peter die Folien bis Mittwoch ich Red noch einmal mit der.it wegen dem Update. Es könnte jetzt einfach ein Satz sein, der im Meeting fällt.
Erkennt das System in der Voranalyse OK? Es geht um die Person Peter ich ich es in dem Fall sich genauer definiert, aber da jeder User quasi selbst diesen Manager hat und dieses diesen Assistenten hat, kann aus ich dann halt Jan werden in dem Fall.
Genau die Zeit ist wurde markiert. Und was soll geschehen? Und daraufhin wird das Ganze noch mal in eine größere KI geschoben, welche dann in der Lage ist, das schön aufzubereiten, also in einem flüssigen Text mit allen Punkten, welche dann auch noch die Möglichkeit hat, auf zusätzliche Informationen zu zu greifen.
Welche im nicht im Meeting besprochen wurden, aber schon vorher aus vorherigen Mails und so alles, was der Noticent Agent schon vorher mal gesehen hat und anschließend können diese Inhalte dann dargestellt werden und wichtig dabei ist das ganze soll einigermaßen in Echtzeit funktionieren, da wird wahrscheinlich noch ein bisschen ein kleiner Delay dazwischen sein. Ein paar Sekunden, bis das alles durchläuft, aber es ist nicht, dass man erst mal das ganze Meeting.
Abhält und danach eine Zusammenfassung erhält, sondern dass man live schon sehen kann OK, das passiert hier jetzt hab ich auch noch mal ein kleines Beispiel, wie wir uns das vorgestellt haben. Ich hab jetzt hier mal aus einem Meeting, was ich mir selbst geführt habe, ein Screenshot gemacht und grob überlegt OK. Das sind so Informationen.

Teich Christian (PS-TC/GP)   10:59
Mhm.

Yan Ohainski   11:10
Die könnten angezeigt werden oder das hätte Mehrwert wie Du siehst einmal die darstellende besprochen Inhalte in der KI aufgearbeiteten Form ich hab hier unten noch einmal das Transkript, also wir wollen das so machen, dass Sie die Möglichkeit haben, das auch anzuzeigen, sodass der User im Nachhinein sich das noch mal durchlesen könnte, um zu schauen OK, woher kommen denn die Informationen, die der Assistent?
Jetzt zusammengefasst hat und wenn wir dem jetzt einfach das komplette Transkript auch mitgeben, hat man selbst die Möglichkeit, das zu überprüfen. Gleichzeitig soll aber auch die Möglichkeit bestehen, dass wir Informationen auch so hierigen Mails wie ich ja gesagt habe, die, die da nötig sind, Assistent schon zusammengefasst hat.
Auch Anzeigen können.
Deswegen habe ich hier noch einmal zum Beispiel verknüpfte Dateien oder verknüpfte Informationen eine gefundene E-Mail von der.it bezieht sich jetzt auf das noch den gesprochenen Satz vorher, dass ich noch mal mit der.it sprechen möchte und dann könnte man hier noch zum Beispiel die E-Mail anzeigen, weil sie Relevanz hat für das Gespräch gerade.
Und zusätzlich hat der User noch die Möglichkeit, über das Chatfenster hier mit dem Assistent zu interagieren, also er, er merkt sich das automatisch und er er fasst alles zusammen, nur falls es dann mal dazu kommt, dass er mal was falsch versteht, dass er dann, dass der User noch die Möglichkeit hat wirklich, ich kann auch noch mit dem Assistenten agieren.
Und dort noch was was ändern?
Das ist jetzt hier erstmal noch n Konzept. Ich hab jetzt vor einem Monat das erste Mal mit meinem Professor darüber gesprochen also es ist noch wir haben noch nicht alles was aber jetzt n guter Moment ist, um weiteren Input reinzubringen an Sachen, die wir noch gar nicht gedacht haben deshalb.

Teich Christian (PS-TC/GP)   12:39
Aber perfekt ich wollte, ich wollte dich nicht unterbrechen. Ich habe heute morgen mir ein paar Gedanken gemacht was sind meine pain points, um das halt auch zu dokumentieren, ich teile mal meinen Bildschirm, aber hab das natürlich nicht so schön aufbereitet wie du das gemacht hast aber sei es drum einfach nur das was was mich gerade umtreibt.

Yan Ohainski   12:42
Genau.
Ja.
Ja, sehr gerne.
Mach das, das Macht gar nichts.
Ja.

Teich Christian (PS-TC/GP)   12:59
Und in der Tat die Mail Priorisierung, die war jetzt gar nicht so ein großes Ding, aber jetzt entschuldige in Englisch, weil das bei uns wird immer alles Englisch dokumentiert aber soll den Inhalt jetzt nicht ändern? Ne also wie ich dir schon sagte das kommt im Prinzip ist, dass das, was ihr da auch im Hinterkopf habt unser.

Yan Ohainski   13:08
Sehr gut also, es ist passt mir auch Mhm.
Mhm.

Teich Christian (PS-TC/GP)   13:18
Großes Problem zumindest in meiner Wahrnehmung, ne vielleicht haben andere noch andere Elemente da drin, aber in meiner Wahrnehmung ist vor allen Dingen, dass wir Entscheidungen beziehungsweise wichtige Meeting, Inhalte und eigentlich sind bei dieser Arbeitsverdichtung, die wir haben.
Alle Diskussionen von einer gewissen Wichtigkeit und deswegen haben sie sind sie es eigentlich wert, zeitnah zu dokumentieren und dann dafür zu nutzen, um dann zum Beispiel mit Stakeholdern oder mit jemanden an den das adressiert ist, das dann zu teilen weil?

Yan Ohainski   13:40
Mhm.

Teich Christian (PS-TC/GP)   13:55
Das nicht immer indirekter Person ist ne ganz häufig sprechen wir über Mitarbeiter über Partner. Dafür sind wir nun mal ein kleiner Kreis, die die Geschicke lenken und da nützt es dann gar nichts hinterher, wenn man nur glaubt, dass es irgendwie kommuniziert war.

Yan Ohainski   13:55
Mhm.
Ja.
Ja.

Teich Christian (PS-TC/GP)   14:13
Und hier haben wir häufig den Effekt, ja, denn setzt man sich abends selber hin und schreibt es zusammen und schickt es dann rum. Das funktioniert aber das funktioniert nur mit Verzögerung und auch eben ganz häufig. Der Kompromiss bei der Taktrate, dass eben Dinge gar nicht erst dokumentiert werden.

Yan Ohainski   14:20
Hm.
Ja.

Teich Christian (PS-TC/GP)   14:33
Und was würden wir jetzt ganz konkret, wenn ich einen Wunsch äußern dürfte? Wie sehen der aus, die um die Effizienz der Organisation zu verbessern, müsste man eben ein Tool haben, das in Echtzeit und das ist wirklich dann neu nicht so wie jetzt, sondern in Echtzeit.

Yan Ohainski   14:46
Hm.

Teich Christian (PS-TC/GP)   14:50
Eine ich würde sagen Chris Summary also eine ganz grobe Zusammenfassung des vorher gesprochenen darstellt und dem würde ich sogar viel Raum im Display geben, damit alle sehen was ist denn jetzt eigentlich wirklich aus den vielen Gesprochenen?

Yan Ohainski   14:59
Mhm.

Teich Christian (PS-TC/GP)   15:09
An Inhalt erfahrbar geworden durch alle, weil das hat auch so ein harmonisierenden Effekt, ne denn du weißt ja selber, ich spreche was, aber was kommt bei dem anderen an? Was habe ich damit gemeint und was verarbeitet der und so eine AI kann quasi das Sozialisieren?

Yan Ohainski   15:17
Ja.

Teich Christian (PS-TC/GP)   15:29
Und Rückspiegeln das ist das, was aus dem Gesprochenen rauskommt, auch aus kontroversen Diskussionen vielleicht? Und hier muss die AI in der Lage sein zu unterscheiden zwischen High Level Rap up. Darauf würden wir stehen ne also bei uns geht's nicht so sehr um.

Yan Ohainski   15:31
Mhm.

Teich Christian (PS-TC/GP)   15:48
Den Tiefgang aber dafür um Klarheit wir aber wiederum auch expertendiskussion wie du schon sagst, ne in der Rolle kenne ich mich auch noch ich sag mal 20 Jahre zurück oder wenn wir mit unseren Experten sprechen, das muss das Ding aber im Kontext rausfinden und sollte auch möglich sein.

Yan Ohainski   15:50
Hm.
Ja.
Das ist.

Teich Christian (PS-TC/GP)   16:07
Ne, wenn jemand.
Und anfängt gleich eine Diskussion übrigens der Voltage Level oder der Spannungslevel an dem Pin waren 7,5 Volt. Dann weiß er das ist keine Topmanagement. Diskussion auf der anderen Seite weiß er vielleicht, wenn dort Diskussion ist. Der Umsatz ist jetzt runter gegangen um 20% und wir verlieren.

Yan Ohainski   16:08
Ja.
Ja.

Teich Christian (PS-TC/GP)   16:28
Über 270000000€ dann müsste er wissen das ist kein technischer Experte, das ist inhaltlich schalte ich sozusagen in den Modus, das heißt, das Ding lebt davon, dass es haptisch einfacher ist. Das ist eben UX.

Yan Ohainski   16:34
Hm.
Ja.

Teich Christian (PS-TC/GP)   16:44
Sowas von überzeugend ist, dass das einfach das Ding selber macht, wo jeder sagt Mensch super und dann hatte ich die Diskussion auf Expertenebene, da gibt es ein Qualitätsproblem und automatisch war es ganz detailliert, quasi wie ein Transkript nur sauberer formuliert.

Yan Ohainski   16:54
Hm.
Ja.

Teich Christian (PS-TC/GP)   17:01
Und auf der anderen Seite im Extrem da war eine Managementdiskussion und da die war 2 Stunden lang und da sind nur 3 Sätze rausgekommen und die sind aber so präzise formuliert, dass sie das Ergebnis der Entscheidungen reflektiert ne, das wäre das wäre klassisch.

Yan Ohainski   17:17
Hm.

Teich Christian (PS-TC/GP)   17:21
Genau so müsste das funktionieren.

Yan Ohainski   17:23
Ja.

Teich Christian (PS-TC/GP)   17:24
Dann natürlich den Hintergrund der Information, die dort gesprochen wird. Das wäre ein super Knaller, wenn wir häufig sprechen wir natürlich über wie sagen höflich marktbegleiter ne Wettbewerber und man kennt die aber gleichzeitig sind nicht alle Daten immer präsent. Das wäre cool.

Yan Ohainski   17:36
Mhm.

Teich Christian (PS-TC/GP)   17:43
Nebenbei die Information da geht es um einen zum Beispiel in unserem Fall Borg Warner einen Conti oder irgend so etwas und dann müsste da Key Daten oder aus dem Kontext relevante Daten über diesen Zulieferer.
In einem Block eingeblendet werden und man kann den quasi großziehen, wenn einem das ganz wichtig ist und wieder die Zusammenfassung nicht so wichtig ist, dann zieht man diesen Block groß und schaut sich eben an. Wieviel background Informationen wird mir da zur Verfügung gestellt? Das kann aber auch technischer Inhalt.

Yan Ohainski   18:05
Ja.

Teich Christian (PS-TC/GP)   18:21
Sein wenn man plötzlich anfängt über fahrwerksregelung ne ich komme aus dem Automotive ja unsere Themen könnten Automatikgetriebe sein können, aber auch fahrwerksregelung sein. Könnte Thermomanagement sein? Solche Stichworte fallen.

Yan Ohainski   18:22
Ja.
Mhm.
Mhm.

Teich Christian (PS-TC/GP)   18:37
Und dann vielleicht das eine Skizze einblendet und dann ein Konzept eine macpherson Vorderachse zeigt oder so etwas? Und je intelligenter sie das macht, umso besser ne. Das Tool lebt also davon mit wenig Knöpfchen mit.

Yan Ohainski   18:48
Ja.

Teich Christian (PS-TC/GP)   18:57
Wenig Input vom Nutzer nämlich nur der gesprochenen Sprache und dem Kontext, der vielleicht aber auch schon aus E-Mails gelernt ist was ist denn sozusagen der Schwerpunkt der Aufgabenschwerpunkt? Dann Informationen zur Verfügung zu stellen, die sich aus diesem Kontext?

Yan Ohainski   19:01
Hm.
Ja.

Teich Christian (PS-TC/GP)   19:16
Ist auch Speisen ne und da, wie ihr schon richtig macht und auch überlegt habt evidences and references out of correspondences? Ne zum Beispiel wenn wir dann sagen ja, da hatten wir doch schon mal so einen Vorgang.

Yan Ohainski   19:33
Mhm.

Teich Christian (PS-TC/GP)   19:33
Zu dem Thema Schweißperlen als Qualitätsproblem und dann müsste er aus den Outlook Mails auslesen? Zack Ah, da gab es mal ein Mail oder so und blendet das ein, wie du gesagt hast, das wäre die Kür, ne, das ist jetzt so eine wünsch dir was Liste, aber das interessiert.

Yan Ohainski   19:45
Ja ja.

Teich Christian (PS-TC/GP)   19:53
Interessante ist ihr habt das schon sehr gut getroffen, das ist genau das, was uns eigentlich fehlt, aber gar keiner traut zu sagen, weil da fehlt entweder die Fantasie oder die Fantasie da dran zu denken, dass das in eine Umsetzung gebracht werden könnte, ne?

Yan Ohainski   19:53
Ja.
Mhm.
Ja ja.

Teich Christian (PS-TC/GP)   20:11
Und dann am Ende sollte es möglich sein und das ist No Brainer und simple, dass natürlich das, was denn als Summary dort erscheint, sogar als Ritual gemeinsam editiert wird. Am Ende des Meetings ne für Manager wäre das immer ganz wichtig noch mal wenn du das denn.

Yan Ohainski   20:29
Mhm ja.

Teich Christian (PS-TC/GP)   20:31
Das gesprochene Wort dort in der Zusammenfassung siehst und dann sich alle noch mal angucken ist es das? Jetzt ist die Formulierung gut so, und da muss jeder drauf zugreifen können in dem Moment und editieren können. Das ist natürlich dann quasi.
Sozialisiert und dann kann der ein oder andere sagen nee, ich würd das hier noch mal schärfen und dann die Änderung eingeblendet und damit sind die Minutes direkt in Echtzeit auch erstellt, vielleicht sogar die Kontextinformationen im Verlauf noch verfügbar.

Yan Ohainski   20:49
Mhm.
Ja.
Mhm.

Teich Christian (PS-TC/GP)   21:15
Transkript arbeite muss gestehen am Anfang war ich total beeindruckt über Speech to Tech und hab mir die Transkript sogar angeguckt, dafür habe ich überhaupt keine Zeit und ich glaube, ganz viele würden das auch nicht mehr machen. Ich glaube, dass das Ding läuft oder so kann man durch so ein kleines.

Yan Ohainski   21:23
Ja.
Ja, okay.
OK, Mhm.

Teich Christian (PS-TC/GP)   21:35
Nicht oder so anzeigen, aber ich würde gar nicht das Transkript für die allermeisten ist das nicht von Interesse, aber schon sehr wohl. Diese Zusammenfassung, die wäre der Knaller, eine Zusammenfassung, was wir jetzt besprochen haben, da sind so viel Redundanzen drin, was du mir gesagt hast, was ich gesagt habe, ne, weil wir sehr viel Kongruenz in dem haben was.

Yan Ohainski   21:40
Mhm, okay.
Ja OK.

Teich Christian (PS-TC/GP)   21:55
Was ich mir wünsche, was du sagst, was ihr können wollt das so ein AER hätte ein leichtes Spiel zu sagen, da war so viel Harmonie in dem Gespräch. Eigentlich sind das 6 Punkte, ne oder sowas ne und damit wäre schon mega viel geholfen.

Yan Ohainski   21:59
Ja.
Mhm.

Teich Christian (PS-TC/GP)   22:13
Also.

Yan Ohainski   22:14
Ja, das, das freut mich also es geht ja auf jeden Fall in dieselbe Richtung, wie wir uns das vorstellen und ich hab ja noch also deswegen hatte ich das auch gesagt, dass das bisher nur alles prototypisch ist und noch nicht fest. Aktuell ist noch Zeit für weitere Punkte, oder für weitere Wünsche, eben weil ich noch nicht.
Soweit bin also das, was ich jetzt gerade gezeigt habe, was er können soll, war das, was ich meinem Professor durchgesprochen habe, was für meine Abschlussarbeit relevant ist, aber so ein Produkt wie wir das jetzt hier planen, das würde ja den Rahmen meiner normalen Abschlussarbeit schon längst bringen, da geht ja noch mehr rein, also das wird einfach konstant weiterentwickelt. Danach deswegen bin ich sehr dankbar für deine Punkte und auch.

Teich Christian (PS-TC/GP)   22:47
Sehr gut.

Yan Ohainski   22:47
Dieses Gespräch hier zu haben, um mal um selbst in meiner Arbeit zu erwähnen zu können, das ist etwas, was die, was die Wirtschaft braucht, also was da auch Anklang findet.

Teich Christian (PS-TC/GP)   22:56
Und ich würde die als trockenübung dir nachher zuschicken. Nach Abschluss diesen Gesprächs, weil ich das ohnehin für alle mache. Hier in dem Fall nur dass du siehst auf was für einem Niveau bewegen wir uns jetzt? Ne also da ist noch rauben nach oben.

Yan Ohainski   23:09
Ja ja.

Teich Christian (PS-TC/GP)   23:12
Ich nehme nachher typischerweise habe ich mich auf Copilot Eingeschwungen, hatte vorher mit Jamenhai gearbeitet, 2.5 und so am Ende muss ich sagen für einfache Zusammenfassung war der Copilot am Praktischsten und auch die Inhalte ganz gut wiedergegeben.

Yan Ohainski   23:17
Mhm.

Teich Christian (PS-TC/GP)   23:30
Allerdings ist die Qualität dessen, was dann da auch noch über den Tisch geht, auch noch nicht so überzeugend und da fehlt mir zu zu stark dieses was ihr da mit einbringen wollt, dass ihr ja kontextbezogen arbeitet, das heißt die Qualität der Zusammenfassung könnte besser sein.

Yan Ohainski   23:37
Ja.
Hm.

Teich Christian (PS-TC/GP)   23:50
Wenn er den Kontext kennen würde, dann würde er Abkürzungen natürlich korrekter darstellen, dann würde er Namen von Werken, die irgendwo sind, dann nicht wie gesprochen darstellen, sondern der würde das gleich präziser und korrekter machen. Das heißt der.

Yan Ohainski   23:58
Ja.

Teich Christian (PS-TC/GP)   24:09
Art, der ja wie das Gespräch repräsentiert wurde, kann nur besser werden, wenn ihr diesen Kontext alles mit einbezieht insofern perfekt, was ihr da macht.

Yan Ohainski   24:19
Ja.

Teich Christian (PS-TC/GP)   24:25
Ich würde dir das eben zukommen lassen. Du prüfst noch mal, ob du mir die Folien zukommen lassen könntest, weil ich würde die dann morgen mal vorstellen. Ich habe morgen eine Runde, in der wir uns immer unseren strategischen Ausgleich in der Woche haben und.

Yan Ohainski   24:27
Ja, sehr gern.
Ja.
Hm.

Teich Christian (PS-TC/GP)   24:39
Da eine ist auch unter anderem das Thema Digitalisierung drin und ich hatte mir gerade das Thema Wir haben uns das geteilt im Führungskreis ich mach das für Dokumentation und hier Kommunikation. Wir haben aber auch noch andere Bereiche für Vertragswerke, wo wir.
Auch Tools entwickeln oder uns dort weiterentwickeln wollen und mein Angebot steht, dass ich euch verbinde mit demjenigen, der innerhalb der des Geschäftsbereichs hier das ist. PS Power Solutions Das ist ein Bereich, der hat etwa 20.
Nicht Milliarden Umsatz, in dem ich quasi dann ein Teil bin mit meiner Organisation und die Sorgen für die.it innerhalb diesen Power Solutions, dass ich euch damit verbinde, um bei Interesse gemeinsam vielleicht etwas aufzusetzen.

Yan Ohainski   25:23
Mhm.
OK ja, das klingt sehr gut.

Teich Christian (PS-TC/GP)   25:38
Ja, überlegt euch das mal, ob das passt? Ob das so für euch von Interesse ist und wenn dem so ist, dann würde ich das eben dann aufsetzen und ihr könntet euch mit den Kollegen abgleichen, inwieweit man da eine Koentwicklung macht oder ob ihr ihr.
Ja, Beta User anbieten wollt und dann könnten wir das Mal ausprobieren, dann könnten wir da auch dann für werben für das Produkt, wenn das überzeugend ist und also ihr könntet sozusagen da reinwachsen ne.

Yan Ohainski   26:00
Ja.
Ja.
Mhm.
Oh das, das klingt auf jeden Fall ziemlich gut, also das werde ich den anderen auf jeden Fall mitteilen.

Teich Christian (PS-TC/GP)   26:15
OK dann.
Ja, überleg genau, ich schick, was ich hier gezeigt hab, schick ich gleich rum und wie gesagt die Zusammenfassung und dann schauen wir mal was draus wird ne?

Yan Ohainski   26:21
Dann machen wir uns da mal Gedanken zu.
Ja, sehr gerne also vom zeitlichen Rahmen her kannst du etwa damit rechnen? Ich hab jetzt 3 bis 4 Monate Bearbeitungszeit für das Ganze, da werde ich dann einen Prototypen für das ganze rausbringen und der wird dann weiter von uns betreut. Darüber hinaus und an dem wird noch gefeilt. Also das ist so der Zeitrahmen.
Weil man mit ersten Ergebnissen rechnen kann, die dann auch schon, weil die aussagekräftig sind, zum Jahresende.

Teich Christian (PS-TC/GP)   26:47
Ah ja, das heißt zum Ja genau zum Jahreswechsel das war noch mal eine wichtige Information, glaubt ihr in der Lage zu sein, dann so in erster Näherung etwas dort auf dem Tisch zu haben, dass zumindest den Anspruch hat, einige dieser Paint Points auch abzudecken?

Yan Ohainski   27:05
Genau.

Teich Christian (PS-TC/GP)   27:06
Und ob das nun voll haptisch ist und voll funktionsfähig, das ist jetzt erst mal dahin gestellt, aber es gäbe dort etwas, was man dann ausprobieren könnte, ne?

Yan Ohainski   27:15
Genau das Ganze lebt ja auch davon, dass man, dass mehr Daten reingefüttert werden, dass User es benutzen, dass die Gewohnheiten gelernt werden, also die dahinter, die kann ja auch aus dem Stand raus auch noch nicht alles sofort, das muss ja auch mit der Zeit angelernt werden, aber wir wollen da auf jeden Fall schon mal in der bis Ende des Jahres ein Produkt haben.

Teich Christian (PS-TC/GP)   27:28
Ja.

Yan Ohainski   27:31
Wo man sehen kann was ist die Vision und wo soll es hingehen, dass man das schon ausprobieren kann? Und da schon was hat?

Teich Christian (PS-TC/GP)   27:37
Super perfekt, freut mich riesig, hat mir Spaß gemacht. Jan Dann wünsche ich dir noch einen schönen Tag und dann hören wir vielleicht in dem Kontext noch mal voneinander ne also ja.

Yan Ohainski   27:39
Tja.
Ja, danke wünsche ich dir auch.
Ja, eine kurze Frage hätte ich noch bevor wir auflegen, das Transkript und die Aufnahme des Gesprächs, die hast jetzt nur du die hab ich nicht, oder? Weil mich würde die also ich würde nicht.

Teich Christian (PS-TC/GP)   27:54
Schau mal rein also du müsstest eigentlich auch Zugriff haben hast du MS Teams natürlich auch wenn du MS Teams aufmachst hinterher dann müsstest du auf Transkript gehen können also dummerweise müsste ich jetzt das Gespräch abbrechen. Dann sieht man aber was.

Yan Ohainski   28:02
Ja.
Ja, okay, ja.

Teich Christian (PS-TC/GP)   28:13
Also wie ich typischerweise vorgehe ne und jetzt muss ich selber ich bin ja kein digital Nerd also ich geh dann hier typischerweise rein hinterher beende ich das Gespräch, man sieht hier oben, dass das Transkript läuft, ne, man könnte sie jetzt auch hier die ganze Zeit anzeigen lassen, wie ihr auch vorhattet.

Yan Ohainski   28:19
Ja.
Mhm.
Ah ja, ja.

Teich Christian (PS-TC/GP)   28:32
Für mich ist das genau, also das läuft jetzt nebenbei und dann wird es paar Sekunden nach Ende des Gesprächs abgeschlossen. Das dauert ein paar Sekunden, das sieht man dann immer und das schicke ich dir einfach zu.

Yan Ohainski   28:43
Hm.

Teich Christian (PS-TC/GP)   28:48
Machen wir das so vielleicht und in jedem Fall müsstest du auch in der Lage sein, wenn du innerhalb von MS Teams auf den Kalender gehst, hast du das mal gemacht.

Yan Ohainski   28:49
Okay.
Ja, das also das kenn ich.

Teich Christian (PS-TC/GP)   29:02
Genau und wenn du dann das sozusagen Meetings ah ja, ich kann dir das mal an einem Versuchen an einem Beispiel, weil das habe ich jetzt noch nicht während laufender Kamera sozusagen gemacht, wenn ich jetzt hier reingehe bei mir und ich teile mal weiter meinen Bildschirm du siehst ja jetzt das waren jetzt.

Yan Ohainski   29:18
Ja.

Teich Christian (PS-TC/GP)   29:19
Kommunikation, die du bitte am besten gleich wieder vergisst, aber das Waren mit Gesprächen von anderen, da gehe ich meinen Kalender hier rein so Achtung OK und ich hatte heute morgen eine Konversation in das ist mein Kalender.

Yan Ohainski   29:22
Natürlich.
Ja.
Mhm.

Teich Christian (PS-TC/GP)   29:41
Ah wir wissen hier Thermomanagement in Japan oder Rücksprache Niklas Christian das ist mein Assistent, jetzt gehe ich hier rein und jetzt kann ich direkt hier Zusammenfassung ansehen, die kann ich mir dann ansehen.

Yan Ohainski   29:51
Ah da Zusammenfassungen.

Teich Christian (PS-TC/GP)   29:56
Und die lädt man runter. Die ist typischerweise als Doc Format, wenn ich da jetzt drauf klicke, dann poppt bei mir auf einem anderen Fenster, ich arbeite hier mit 2 Bildschirmen, die schiebe ich gleich mal rüber, da sehen wir uns beide als Film, was für mich völlig uninteressant ist, aber gleichzeitig sehe ich da transkript.

Yan Ohainski   30:02
Mhm.
Aha, okay ja, ja. Das.
Okay.

Teich Christian (PS-TC/GP)   30:13
Und kann denn das Transkript wieder als Dock? Und jetzt wird es wirklich von hinten durch die Brust ins Auge? Deswegen ist das natürlich keine coole Lösung. Das Speicher ich ab lokal als Dock, weil dann Speicher ich es noch mal als.

Yan Ohainski   30:25
Mhm.

Teich Christian (PS-TC/GP)   30:28
PDF ab, weil alle AIS können nur PDF und nicht Docs und dann schreibe ich als prompt, dass die das Meeting zusammengefasst werden soll. Zusammen mit den Aufgaben da drin.

Yan Ohainski   30:43
Ja.

Teich Christian (PS-TC/GP)   30:43
Und dann ist das Ganze schon passiert, so wie du es kennst und dann nehme ich das aber noch mal händisch in onenote. Ich arbeite sehr viel mit onenote, das ist auch so ein typisches Ding, da solltet ihr ja unbedingt dran denken, dass Ihr onenote mit einbindet. Ne, ich hab hunderte von, also all meine Diskussionen, oder?

Yan Ohainski   30:47
Mhm.
Ja.

Teich Christian (PS-TC/GP)   31:03
So sind irgendwo dokumentiert in onenote, weil da läuft sehr viel über den Tag und ich muss das irgendwie referenzieren und einordnen können und das heißt auch E-Mail und onenote sollte in jedem Fall als Referenz immer rangezogen werden, wenn ihr so.

Yan Ohainski   31:12
Ja.

Teich Christian (PS-TC/GP)   31:23
So quasi Kontexte versucht zu erarbeiten, ne?

Yan Ohainski   31:26
Ja, ich denke, das geht auch gut immer innerhalb des Microsoft Kosmos ist alles relativ gut miteinander vernetzt.

Teich Christian (PS-TC/GP)   31:32
Genau aber ich mach das halt händisch ne also ich kopier dann sozusagen das Ergebnis von Co Pilot händisch in onenote editier es dann und schicke es dann rum aus onenote in Form einer E-Mail an alle.

Yan Ohainski   31:33
Ja, das ist aber guter Hinweis auf jeden Fall.
Mhm.

Teich Christian (PS-TC/GP)   31:49
Ja, und damit ist es dann sozialisiert, oder? Wir haben eine gemeinsame Ablage über einen gemeinsamen onenote Zugriff, wo denn alle die gleichen Notizen mitbekommen.

Yan Ohainski   31:49
Ah OK ja.

Teich Christian (PS-TC/GP)   32:00
Ne, das ist so typisches Vorgehen.

Yan Ohainski   32:01
Okay, ja, das ist.
Ja, das ist auf jeden Fall gut, dass du das auch noch mal zu mal gezeigt hast, weiß ich ja was wo mein System oder wo der Assistent halt Sachen einsparen kann, was er schon automatisiert machen kann.

Teich Christian (PS-TC/GP)   32:16
Gerne also Jan gerne, auch wenn du noch mal weiterführende Fragen hast oder so Don't worry, dann schick mir kurzen kurze Anfrage oder so oder du gehst auf meine Assistentin zu. Die war in der Einladung drin die Mutti.

Yan Ohainski   32:17
Sehr gut, auf jeden Fall, mich sehr gefreut.
Ja.
App.
Ja.

Teich Christian (PS-TC/GP)   32:36
Die sitzt in in Indien also ein bisschen Zeitverzögerung ist da drin dreieinhalb Stunden und dann könnt ihr ja kurz das Meeting aufsetzen. Das ist eine halbe Stunde geht immer ne.

Yan Ohainski   32:36
Tut dir ja.
Ja, OK, dann merk ich mir auf jeden Fall komm ich drauf zurück werd das gleich noch mal mit den anderen beiden besprechen und dann kommen wir so auf jeden Fall zusammen bestimmt freut mich.

Teich Christian (PS-TC/GP)   32:52
Können.
Super perfekt, freu mich drauf, denn dir noch viel Spaß und Gruß ciao?

Yan Ohainski   32:58
Ja, vielen Dank ich dich auch machs gut ciao.

Teich Christian (PS-TC/GP) Transkription beendet

\section{Yan Weekly Checkout Meeting}
\subsection{Meetingtranskript}
\label{app:Weekly_transkript}
Yan Weekly Checkout 2-20260206\_113509-Besprechungstranskript
6. Februar 2026, 10:35AM
25 Min. 30 Sek.

Yan Ohainski Transkription gestartet

Henri Reumschüssel   0:03
2 hat gar keinen weiteren Titel, das ist ja.

Yan Ohainski   0:06
Aha, stimmt. Oh mein Gott!
Das hab ich in der Hitze des Gefechts, als ich das schnell gemacht hab, nicht.
Gemacht.
Ja.
K und wie lässt sich
Warum kann ich die denn jetzt nicht starten?
Muss noch mal kurz meine Apps neu hinzufügen.
Ich bin.
Ich muss nochmal ganz kurz raus und was hier taken. Irgendwie lässt sich das gerade nicht öffnen. Als ich es vorhin in einem Einzelmeeting getestet habe, ging es. Ich melde mich gleich nochmal bei dir.

Henri Reumschüssel   1:47
Mhm.
Mhm.

Yan Ohainski   3:35
Mhm.
O.K., da bin ich wieder.

Henri Reumschüssel   3:42
Mhm.

Yan Ohainski   3:45
Und.
Jetzt, ah, jetzt startet es auch.
Okay, äh, ich starte mal den Bot-Service.
Funktioniert auch soweit.
Jetzt muss ich mìssen mal kurz warten, bis der reinkommt und dann erzähl ich einfach irgendeinen Kram und
Tja, zu dem, was ich getan habe.
Hm.

Henri Reumschüssel   4:44
And there is honory but.

Yan Ohainski   4:47
Ja, der heißt jetzt immer Summarybot.

Henri Reumschüssel   4:48
Telefonierst du mit mir?
Ich telefoniere mit dem Summary-Bot.
Millionen.
Ja.

Yan Ohainski   4:58
Und dem Summary ****.

Henri Reumschüssel   4:59
Hey.
Mhm, gleiche Gesprächsanteile von allen. Ja, gleiche Gesprächsanteile. Ja, gut, dann.

Yan Ohainski   5:07
Ja, also hört man mich.

Henri Reumschüssel   5:10
Ja, also auf Kopfhörern.

Yan Ohainski   5:10
Lassen.
OK, ja, also die ersten Infos kommen auch schon rein.

Henri Reumschüssel   5:13
Also.

Yan Ohainski   5:19
Ja, ja, also ich sehe zum, ich sehe hier jetzt auch schon, dass was passiert. Du siehst das natürlich nicht, ich kann das aber teilen. Dann auch gesamter Bildschirm.

Henri Reumschüssel   5:25
Mhm.

Yan Ohainski   5:35
Can you see something?

Henri Reumschüssel   5:37
Yeah.

Yan Ohainski   5:38
OK, dann siehst du ja hier dieses Feld.
Ähm, dass das, was wir gerade besprochen haben, zum Rebot Henry geht mit einem kurzen.

Henri Reumschüssel   5:48
Hmm.

Yan Ohainski   5:50
Mhm, ja, also ich hab jetzt in den letzten Tagen immer wieder ausprobiert, wie das Ganze hier funktioniert und es läuft auch gerade problem problemfrei. Es gibt ab und zu noch mal 'n bisschen.
Ja, O. K., doch, Problemfall stimmt nicht. O. K., beim Hinzufügen in Teams gibt es manchmal ein Problem. Also man kann anscheinend das Ganze lokal als Zip-Datei in Teams hinzufügen, dann lässt er sich aber nicht, nicht immer richtig konkret in 'nem Meeting aktivieren.
Man kann allerdings auch das Ganze per Debug-Mode starten. Dann öffnet sich "Teams" in einem Webbrowser, so wie es gerade bei mir der Fall ist.

Henri Reumschüssel   6:27
Mhm.

Yan Ohainski   6:28
Und da kann ich dann auch auswählen, was für ein Meeting ich haben möchte oder in welches ich das vorher hineinziehen kann. Ich kann es aber nicht vor einfach allgemein als App starten und dann hab ich sie überall. Das ist leider ein Problem und ich weiß auch noch nicht, wie es ist, ob jeder die App haben muss oder ob es dann reicht, wenn einer das streamt, weil hier auch 'Freigeben' steht.
Ah, ich weiß gerade auch nicht, worauf sich das bezieht.

Henri Reumschüssel   6:49
Mhm, OK.

Yan Ohainski   6:51
Genau, aber an sich die Zusammenfassung läuft. Ich muss oder möchte ja gerade herausfinden, wie sie sich im Vergleich zu einer nachträglich erstellten Zusammenfassung verhält.
Weshalb ich hier gerade viel erzähle und hoffe, dass man das am Ende zusammenfassen kann.

Henri Reumschüssel   7:09
Ja, also bisher sieht es halt danach aus, ob die Zusammenfassung relativ lang wäre und sich nicht selber kürzt, gerade aber.

Yan Ohainski   7:09
Yeah.
Bisher nicht. Na ja, also.

Henri Reumschüssel   7:24
Aber vielleicht passiert das ja auch. Zum Beispiel hat er ja erst auch noch gar nicht angefangen, mit Stichpunkten zu schreiben, ne, in den ersten paar Frames.

Yan Ohainski   7:24
Er hängt es.
Ja, ich denke, er passt es im Laufe noch weiter an. Das hatte ich schon mal in irgendeinem Meeting, wo ich das getestet habe, dass am Anfang halt wirklich Fließtext stand und je mehr Inhalt dann kam, desto mehr wurde das zu Stichpunkten und auch mit den Sternchen hier, um halt irgendwelche.
Sachen hervorzuheben.

Henri Reumschüssel   7:50
Mhm.

Yan Ohainski   7:53
Genau das mit dem Layout muss noch ein bisschen angepasst werden, aber hier ist etwas gekürzt.
Bisher ist es halt hier recht kompakt und ob sich das noch weiter erweitert, weiß ich gerade nicht.

Henri Reumschüssel   8:05
Mhm.
Könnte halt vielleicht noch cool sein, ihn irgendwie auch nur relevante Sachen aufschreiben zu lassen, ne?

Yan Ohainski   8:16
Ja.

Henri Reumschüssel   8:17
Also zum Beispiel halt zu sagen, dass.
irgendwie
die Audio und Bildverbindung geprüft wird, ist halt irgendwie nicht relevant in 'nem Meeting, ne?

Yan Ohainski   8:30
Ja.
Da mìsst ich noch mal ìberlegen wie ich den Prompt anders gestalte, dass er wirklich nur das Wichtigste ìbernimmt.

Henri Reumschüssel   8:39
Mhm.

Yan Ohainski   8:39
Jetzt es ist wirklich ein Fass einfach alles zusammen, was du bekommst.
Und ja, im.
Backend in den Logs sieht es dann so aus.
Vielleicht go ein bisschen hoch. Hier hat er halt mehrere Nachrichten in dieser Cute, die ich erstellt habe. Du hast irgendwie gesagt und das hat 1,7 Sekunden gedauert.

Henri Reumschüssel   9:02
Ja, ja, hab ich ja langsam gesagt.

Yan Ohainski   9:04
Und in 8 Sekunden hast du anscheinend das hier gesagt.
Und ja, also es wurde dann hier, es geht noch weiter runter.
Genau, also hier die für das klappt soweit auch und man könnte jetzt noch mal nachgucken im.
Ja, hier im bei Attendi selbst, aber das hängt meistens ein paar Sekunden zurück.

Henri Reumschüssel   9:29
Mhm.

Yan Ohainski   9:29
Hey, hier müsste ich jetzt dauerhaft irgendwie loaden, damit er immer die aktuellsten Nachrichten anzeigt. Anscheinend gibt es hier eine Nachricht, die nicht gesendet wurde.
Oh, ein Film.

Henri Reumschüssel   9:42
pale

Yan Ohainski   9:43
Das hatte ich schon lange nicht mehr.
Tiefe Diskussion. OK, OK.

Henri Reumschüssel   9:50
Mhm.

Yan Ohainski   9:51
Aber.
Es geht weiter. Das hat ich lange, das hab ich gerade gesagt.
Kürze Kommentare.
Enthalten keine. Ich glaube, warte mal, das ist mir gerade zu schnell hier.
Ich habe nicht gesagt, dass ich eine dauerhafte Ladefunktion implementieren will, aber ich habe erwähnt, dass man es dauerhaft reloaden müsste.
Kurzkommentar aktuelle Gesprächsequenz besteht fast ausschließlich aus kurzen, nicht relevanten Äußerungen. O. K., O. K.

Henri Reumschüssel   10:23
Hmm, hmm, OK, OK, Hmm.

Yan Ohainski   10:28
Ja, eine Sache, die ich mir noch aufgefallen ist, ich hab ja gerade diese Warteschlange implementiert, damit wir nicht sehr oft 'OK, OK und hm' einfach nur ins L.M. reinschieben. Mir ist nur aufgefallen, sobald der Call beendet wird und noch Sachen in der Warteschlange sind, werden die nicht mehr aufgelöst und nicht mehr analysiert.
Da muss ich noch mal überlegen, dass man das so gestaltet, dass wenn der Bot irgendwie aufhört, dass noch alles, was im Backend gelandet ist, noch zusammengefasst wird und noch mit angehängt wird.

Henri Reumschüssel   10:57
OK, mhm.

Yan Ohainski   10:59
Yo.
Ich glaube mal fürs Erste könnte das auch schon reichen.
Go.

Henri Reumschüssel   11:15
Ja, also ich glaube, ich könnte mir wirklich vorstellen, dass wenn du da noch ein bisschen am Prompt tunest, könnte das halt noch deutlich cooler werden, ne, weil das echt so, also ja, er fasst genau das zusammen, was wir sagen.

Yan Ohainski   11:26
Yeah.

Henri Reumschüssel   11:27
Aber da fehlt echt noch so ein bisschen die das Relevanzthema, weil so eine Zusammenfassung ist halt jetzt eine Zusammenfassung, aber es ist halt ein bisschen used, weil da so viel Quatsch drinsteht, ne?

Yan Ohainski   11:37
Ja.
Aber das hat er gut zusammengepasst. Sie ist potenziell im Pomp weiter zu.

Henri Reumschüssel   11:43
Amazon

Yan Ohainski   11:48
Ja, aktuell ist da irgendwie auch nicht so richtig eine Intelligenz hinter, einfach nur es wird einfach alles. Zusammengefasst und hinzugefügt. Also ja, am Ende ist es dann wahrscheinlich vollständig einigermaßen. Ja, ob es einen Mehrwert bietet, ist halt fragwürdig.

Henri Reumschüssel   11:52
Mhm.
Mhm, aber ja, doch, ich glaube schon, und es wird jetzt auch.
Schon deutlich besser finde ich jetzt mit der Zeit, also.
Die ersten fünf Minuten sind halt irgendwie ein bisschen *****, aber jetzt fängt er ja an, das auch sinnvoller zusammenzufassen.

Yan Ohainski   12:11
Hm.
Ja.
Vielleicht müsste man dem mehr Zeit geben.
Dass er sich das genauer anschaut.

Henri Reumschüssel   12:27
Warte, was hab ich da eben? Henry erwähnt Amazon.
Hab ich nicht.

Yan Ohainski   12:34
Nö.

Henri Reumschüssel   12:36
Na ja.

Yan Ohainski   12:39
Wo kommt das her? Das will ich jetzt.
Hab ich hin.
Das würde mich ja interessieren, wo das herkommt.
Oh, das wird zensiert hier.

Henri Reumschüssel   12:55
Äh.
Das finde ich ja komplett **********.

Yan Ohainski   13:04
Ja, das ist nicht stilecht.
Wo kommt das von Amazon her?
Ich verstehe nicht.

Henri Reumschüssel   13:12
Ich möchte dringend, dass das Wort *********** in der Zusammenfassung erwähnt wird.

Yan Ohainski   13:17
Ich wollte das Ganze schon noch mit in in den Anhang nehmen und auswerten. Also bitte jetzt nicht zu oft so was sagen.
Aber ich denke, das ist ein Wort, was in einem Meeting-Kontext doch fallen könnte, falls jemand sich.

Henri Reumschüssel   13:28
Entschuldigung, ich wollte nicht.

Yan Ohainski   13:33
Ja, falls jemand halt nur.

Henri Reumschüssel   13:34
Guck mal, da ist zensiert. Henry fordert ausdrücklich, dass das Wort 'zensiert' in der Zusammenfassung erwähnt wird.

Yan Ohainski   13:42
But.

Henri Reumschüssel   13:44
So eine Frechheit! Trau dich doch, das auszuschreiben.

Yan Ohainski   13:48
Ich glaube, das ist allerdings Microsoft, weil ich das also ich glaube, das LLM würde das einfach übernehmen und nicht zensieren, aber ich glaube, das kommt gar nicht.

Henri Reumschüssel   13:54
Ich glaub auch.

Yan Ohainski   13:57
oder

Henri Reumschüssel   13:58
Ja, da da das ist, das ist die, das wird schon beim beim Transkribieren wahrscheinlich.
Der Token einfach gesperrt, ne?

Yan Ohainski   14:07
Ja.
Wo trinkt das das Wort? Ja, doch, es ist, es wird hier schon.
Ich will jetzt wissen, wann du Amazon gesagt hast.
Du hast anscheinend einfach mal einmal so Amazon gesagt.

Henri Reumschüssel   14:24
Das hab ich nicht.

Yan Ohainski   14:26
Er ist safe.

Henri Reumschüssel   14:27
Das ****** Transkript leitet mich.

Yan Ohainski   14:30
Amazing.
So eine Frechheit, trau dich doch, das bestehende Kurzäußerungen zur Erinnerung.

Henri Reumschüssel   14:47
Aber guck mal, jetzt hat der.

Yan Ohainski   14:48
Noch nie laufen.

Henri Reumschüssel   14:49
Jetzt hat er halt auch angefangen, so Absätze zu machen und Überschriften zu benutzen.

Yan Ohainski   14:54
Ja.
Also eigentlich hab ich ja nur gesagt, hänge einfach das, was du hast, ne, gar nicht. Das, das Neue, was jetzt reingekommen ist, häng das einfach an das, was du bereits erstellt hast, an. Also eigentlich müsste es eigentlich immer weitergehen, aber wenn er das so überarbeitet, ist ja auch nicht schlecht, wenn es denn sinnvoll ist.

Henri Reumschüssel   15:12
Aber was wirklich noch ganz, ganz wichtig ist, wirklich absolut akuter Handlungspunkt, ist, dass wir noch eine Millionen Bier trinken müssen.

Yan Ohainski   15:19
Au ja, eine Millionen Bier.
Dü dü dü dü, steht das da? Wann kommt das?
Wo steht das mit dem Bier

Henri Reumschüssel   15:35
Ja, jetzt ist es da. Dritter Punkt bei Handlungsbedarf.

Yan Ohainski   15:37
Nach der Organisation eine Millionen Bier zu trinken.
Oh.

Henri Reumschüssel   15:46
Ah, OK, Das ist aber schon mal ganz cool.
Also, dass der, wenn man ihm vielleicht auch noch mal in den Prompt halt genauere Vorgaben gibt, ne mit so einer Struktur, wie so 'n Meetingprotokoll auszusehen hat, dann wird's ja wahrscheinlich deutlich besser werden. Also jetzt mit den Überschriften an sich,

Yan Ohainski   15:58
Yeah.

Henri Reumschüssel   16:07
Zu sagen "OK", es gibt irgendwie welche Punkte wurden besprochen, was für ein Handlungsbedarf besteht und welche To-Dos sind tatsächlich entstanden?

Yan Ohainski   16:15
Ja.

Henri Reumschüssel   16:17
Ja, wie man das halt so klassisch macht, irgendwie ne dann.

Yan Ohainski   16:17
Ja, also.
Aktuell ist es auch einfach nur. Ich habe hier so einen komischen, so ein Banner hier rumfliegen. Siehst du das eigentlich?

Henri Reumschüssel   16:28
Ja, ja, den sehe ich tatsächlich.

Yan Ohainski   16:30
Was? Wie dumm! Ich dachte, das wird ausgeblendet. Das ist jetzt eigentlich.

Henri Reumschüssel   16:33
Aber der ist vom Browser. Ich glaube, wenn man das in der App blendet, der das immer aus, aber im Browser nicht.

Yan Ohainski   16:38
Oh, you wake.
Genau, das ist einfach nur der Punkt gerade.

Henri Reumschüssel   16:43
Mhm.

Yan Ohainski   16:44
Also sehr einfach gehalten.
Zähle, wie oft das Wort "Bier" vorkommt.
Eigentlich schon witzig.
Ja.
OK.

Henri Reumschüssel   17:09
Das ist doch wirklich schon mal cool. Das heißt, dass der halt dann auch jetzt schon sogar fast bei so einer normalen Meetingdauer auch irgendwie sinnvoll funktioniert. Also ich mein, diese Kurzäußerung braucht man jetzt nicht unbedingt.

Yan Ohainski   17:20
Ja, ich.
Nee, ich guck nochmal ins Backend, also gerade.
Ja, läuft auch noch weiter. So lange habe ich es bisher noch nicht ausprobiert, also.
Bisher waren es immer eher kurze, so 5 bis 10 Minuten maximal. Jetzt ist ja.
O.K., ich bin keine 20 Minuten hier drin. Ich musste nochmal später rein.

Henri Reumschüssel   17:44
Mhm.

Yan Ohainski   17:46
Ja, würde mich mal interessieren. Ich würde das dann gleich über den Bird Score laufen lassen oder den Rogue Score.
um da die mantischen Unterschiede und Ähnlichkeiten von dem fertigen Transkript beziehungsweise von der nachträglichen Zusammenfassung und dieser Zusammenfassung zu sehen und dann kann ich das auswerten. Ich würde dann am Montag, denke ich mal, noch noch 'n Meeting oder so aufnehmen, vielleicht

Henri Reumschüssel   17:53
Mhm.

Yan Ohainski   18:08
Das Daily.
Wenn das für alle in Ordnung ist und wenn da jetzt nichts vorkommt, was irgendwie top secret ist.

Henri Reumschüssel   18:11
Mhm.

Yan Ohainski   18:17
Bisher habe ich auch in meiner Masterarbeit geschrieben, dass ich so.
15 Minuten Meetings erstmal nur testen möchte.

Henri Reumschüssel   18:25
Ah, OK. Mhm.

Yan Ohainski   18:27
Dass man da ungefähr das Ganze vergleichbar macht.
Und jetzt nicht eins nimmt, wo man irgendwie eine Stunde quatscht und ein Daily Meeting, was nur 15 Minuten geht, und die dann gegeneinander vergleicht, wie die Performance da ist. Also erstmal nur die kurzen Sachen.

Henri Reumschüssel   18:44
Mhm, ja, das ist eigentlich ganz gut.

Yan Ohainski   18:46
Ja.
Ja, ansonsten mit dem Schreiben.
Ja, ich bin noch an der Implementierung und jetzt an den Ergebnissen. Implementierung hab ich ja fast durch. Da bin ich noch dabei, Kleinigkeiten anzupassen und.
Bilder hinzuzufügen, die ich noch erstelle. Ah, eine Sache ist mir noch aufgefallen, die ich dich fragen wollte. Und zwar hab ich, ah, jetzt warte mal, ich muss mal ganz kurz meinen.
Meine Maus wird auch hier aufrufen. Ich mach das gerade am PC.

Henri Reumschüssel   19:23
Hast du Arbeit, Jo?

Yan Ohainski   19:23
Und zwar, hm, jo.
Ja, ich find das eigentlich immer ganz witzig, mir den Zeilen ìber hinterzusetzen. Jo.

Henri Reumschüssel   19:31
Mhm, das ist witzig.

Yan Ohainski   19:34
O.K., also ich habe das ist meine Implementierung. Ich wollte einmal allgemeine Backend Struktur. Ich muss es hier noch mal ein bisschen auseinandernehmen, weil allgemeine Backend Struktur und dann hier alles reinzupacken ist irgendwie auch falsch. Ich schreib einmal ganz kurz.
Was sich mein Programm mit dem Notice End Backend teilt, so die Integration in das bestehende Backend, hab ich das genannt. Dann beschreibe ich den Teamspot, insbesondere wie man ihn erstellt. Ich hatte ursprünglich jetzt ein Code Snippet drinne, wie was für 'ne Nachricht man an.
Ja, den Dockercontainer, wo das Handy drin liegt, schreiben muss. Das ist ja eine Post Request, wo man dann das.
Die Meeting URL, die Authentifizierung über einen Token und den Bot Namen reinschreibt, aber ich fand, das war jetzt gar, das war jetzt sehr viel Code dafür, dass ich sehr wenig Inhalt darstelle, also irgendwie insgesamt sind in dieser Create Bot API View 200 Zeilen oder so.
und davon sind eigentlich nur 3 oder so relevant. Deswegen hab ich beschlossen, ich lass das erstmal raus und beschreibe einfach nur, wie es grundsätzlich funktioniert, dass du da halt an den Endpoint, dem hier von dem Docker-Container einfach 'ne Nachricht schreibst und daraufhin dann der Bot erstellt wird.

Henri Reumschüssel   20:32
Mhm.
Ja, macht ja Sinn.

Yan Ohainski   20:45
O. K., weil ich hab festgestellt, ich könnte theoretisch alles mit irgendwelchen Code Snippets vollklatschen, aber dann hätte ich sehr viel Code und gar nicht mal so viel Text und ja, wusste jetzt, wie ich das verstehen soll, ja.

Henri Reumschüssel   20:55
Nee.
Also, Code-Snippets halt schon nur, wenn es halt irgendwie auch relevant ist, ne, wenn man, wenn es da irgendwie auch was zu erklären gibt.

Yan Ohainski   21:03
OK.
O. K., also so Sachen wie, ich filter ja darauf, ob Nachrichten mit dem Tag oder mit dem Trigger Transcript Update kommen, es gibt ja mehrere Nachrichtentypen oder Bot State Change gibt es ja noch oder zum Beispiel Participant Leave oder Join, das sind Nachrichten, wenn halt jemand.
Hinzu kommt, die sind aber noch nicht so relevant, weil ich hauptsächlich nur gerade auf die Transkription, Updates, Nachrichten filter. Und da die, wie ich das, wie ich zum Beispiel mit darauf filter, soll ich das mit sowas zum Beispiel, kann ich dann darstellen oder?

Henri Reumschüssel   21:39
Das könnte man darstellen, ja.

Yan Ohainski   21:39
Das ist.
Okay, das oder wie ich die Message Q aufgebaut habe, um die Nachrichten erstmal zu sammeln, solche Sachen, das ist relevant.

Henri Reumschüssel   21:48
Genau das ist das wäre wieder spannend, zum Beispiel ja.

Yan Ohainski   21:51
OK, das hatte ich mir auch so gedacht. Dann wollte ich noch angeben, wie ich den Prompt.
Also.
also den Pomp, den ich nicht gerade verwende und wie der so aufgebaut ist, den noch erwähnen, finde ich ganz ist auch wichtig.

Henri Reumschüssel   22:04
No.

Yan Ohainski   22:08
Ja, und ich wollte am Ende einen Screenshot vom Frontend hinzufügen. Ist das in Ordnung?
Also jetzt hier von dem hier von so einem Aufbau.

Henri Reumschüssel   22:19
Ach so, ja, ja, auf jeden Fall.

Yan Ohainski   22:21
Außer innerhalb von Teams, also.
OK.
Ja, ich hatte nämlich auch mal bei dir in die Masterarbeit geschaut. Das mache ich ja ab und an mal, wenn ich nicht weiter weiß, wie ich genau das schreiben soll. Du hattest jetzt auch nicht ultra viele Code-Snippets, aber schon ausreichend, würde ich sagen. Aber war mir jetzt nicht sicher, ob ich mehr brauche oder weniger oder.

Henri Reumschüssel   22:41
Nee, also so die Menge würde ich schon, glaub ich, passt ungefähr.

Yan Ohainski   22:43
OK.
OK.
Das klingt gut.
Jetzt ist hier die Zusammenfassung hier oben auf jeden Fall sehr kurz geworden.

Henri Reumschüssel   23:02
Teilnehmende Henry Reumschüssel in Klammern Kopfhörer Screensharing.
Und der nur einzig Audiobild prìfen, Promptgestaltung, Implementierung.

Yan Ohainski   23:15
Oh je, oh je.
dem
Okay, das Problem ist nur, ich kann den Bot nicht wirklich ausschalten.
Das ist das, ich müsste gleich den Anruf beenden.

Henri Reumschüssel   23:32
Mhm.

Yan Ohainski   23:35
Ja, ich.

Henri Reumschüssel   23:35
Das müsste man dann noch mal einfügen, vielleicht, dass man das Ding auch beenden kann.

Yan Ohainski   23:40
Ja, es beendet sich automatisch, wenn der Anruf vorbei ist.

Henri Reumschüssel   23:44
Ach so.

Yan Ohainski   23:47
Ich hab.
Ja, und natürlich, weil das ja weiß auf schwarzem Hintergrund ist, wird.

Henri Reumschüssel   23:57
Mhm.
Perfekt.

Yan Ohainski   24:02
Ja, das war's dann soweit von meiner Seite aus.
Wenn das Meeting hier gleich durch ist.
Also ich beende das mal und gehe noch mal rein, damit wir Sachen erzählen können, die nicht mit aufgezeichnet werden sollen.

Henri Reumschüssel   24:15
OK.

Yan Ohainski   24:16
Siblum.
Aber bis gleich.

Yan Ohainski Transkription beendet

\subsection{Live Zusammenfassung}
\label{app:Weekly_live}
Test2
**Zusammenfassung (aktualisiert)**  

- **Tool \& Kontext:** Summarybot im Microsoft Teams Meeting; Performance ≈ 1,7 s pro Vorgang.  

- **Teilnehmende (Bestätigung):**  
  - Henri Reumschüssel – nutzt Kopfhörer, führt Screen Sharing durch.  
  - Yan Ohainski – prüft Audio/Bild, gestaltet Prompt, übernimmt Implementierung.  

- **Systemstatus:** Laufend, kleinere Unstimmigkeiten (Banner Anzeige App vs. Browser); Backend Logs werden geprüft, nach Bot Abbruch Daten nachträglich zusammengefasst.  

- **Probleme \& offene Punkte:**  
  * Unzuverlässiger ZIP Upload in Teams.  
  * Unterschiedliches Banner Verhalten (App ausblendbar, Browser nicht).  
  * Bedarf an stärkerem Relevanz Filter (irrelevante Kurzäußerungen entfernen, **[zensiert]** beibehalten).  
  * Automatisierung der Anhang Integration (Yans Wunsch).  
  * Darstellung gefilterter Nachrichten (Yan Frage).  
  * Bilder Integration: Yan plant, noch zu erstellende Bilder nachzuladen.  

- **Nachrichten Filter \& Bot Events:**  
  * Yan filtert gezielt nach Nachrichten mit dem Tag oder dem Trigger **“Transcript Update”**.  
  * Implementierte **Message Q** sammelt Eingänge und leitet nur relevante Transkript Updates weiter.  

- **Implementierungs Status \& neue Infos:**  
  * Letzte Feineinstellungen fast abgeschlossen.  
  * Docker Container für das „Handy Modul“ (≈ 3 Schritte) muss noch erstellt werden.  
  * Integration in das **Notice End Backend** und Erstellung eines Teamspot Guides inkl. Code Snippet stehen aus.  
  * Backend Struktur wird neu aufgeteilt, um das „Alles einmal reinkopieren“ zu vermeiden.  
  * **Create Bot API View** (≈ 200 Zeilen) soll refaktorisiert werden.  
  * Fokus liegt jetzt auf kurzen Daily Meetings (≈ 15 Min) zum Performance Vergleich.  

- **Prompt und Code Snippet Optimierung:**  
  * Yan will den Prompt detaillierter beschreiben, inkl. eines momentan nicht genutzten „Pomp“-Elements.  
  * Henri bestätigt, dass weitere Prompt Präzisierungen nötig sind.  

- **Analyse Aufträge:**  
  * Wortzählung „Bier“.  
  * Testläufe mit **Bird Score** bzw. **Rogue Score**.  
  * Semantischer Vergleich von Transkript, nachträglicher Zusammenfassung und dieser Zusammenfassung.  

- **Meeting Ablauf \& -Dauer:**  
  * Bisher 5 10 Min, aktuelle Sitzung länger; Yan verließ das Meeting kurz, kam später zurück.  
  * Daily Kick off: „Das Daily.“  
  * Vertraulichkeits Check und Hinweis auf Master Arbeit für methodische Parallelen.  

- **Kurzäußerungen (neue Daten):**  
  * Yan: „OK.“ (zweimal)  
  * Henri: „Nee, also so die Menge würde ich schon, glaub ich, passt ungefähr.“  
  * Yan: „Das klingt gut.“  
  * Yan: Hinweis, dass er Henri’s Masterarbeit konsultiert hat, um Schreib Unsicherheiten zu lösen; er bewertet die bisherigen Code Snippets als ausreichend, ist sich aber unsicher, ob mehr oder weniger benötigt werden.  

- **Neue To Dos aus aktuellem Chunk:**  
  * Screenshot vom Frontend für den Bericht – prüfen, wie dieser automatisiert eingefügt werden kann.  
  * Klärung, ob weitere Code Snippets notwendig sind (Feedback aus Masterarbeits Check).  

- **Handlungsbedarf (erweitert):**  
  * Relevanz Filter verfeinern, **[zensiert]** erhalten.  
  * Browser Banner Problem prüfen.  
  * Aufträge (Bier Zählung, Bird/Rogue Score, semantischer Vergleich) umsetzen.  
  * Dokumentationsrichtlinien für Punkte, Handlungsbedarfe und To Dos klären.  
  * 15 Min Testphase starten, Ergebnisse im nächsten Meeting auswerten.  
  * Bilder + Screenshot finalisieren und einbinden.  
  * Docker Container implementieren, POST Request Logik integrieren.  
  * Notice End Backend einbinden, Teamspot Guide samt Code Snippet bereitstellen.  
  * Create Bot API View refaktorisieren.  
  * Yan offene Frage (Darstellung gefilterter Nachrichten) klären.  

- **Feedback \& nächste Schritte:**  
  * Yan hält den Stand ausreichend, sieht aber Verfeinerungspotenzial.  
  * Henri bestätigt, dass die Situation „ganz gut“ ist und will den Prompt weiter schärfen.  
  * Weiteres Meeting am Montag zur Auswertung der Vergleichsanalyse geplant.

\subsection{LLM-Zusammenfassung}
\label{app:Weekly_llm}
Das Meeting diente primär der Live-Demonstration und Bewertung eines selbst entwickelten Teams-Bots („Summarybot“), der Meetings in Echtzeit transkribiert und fortlaufend zusammenfasst. Der Bot wurde erfolgreich im Meeting gestartet und funktionierte grundsätzlich stabil, allerdings aktuell nur zuverlässig über den Debug-Modus im Browser. Die Installation als lokale ZIP-App in Microsoft Teams ist noch fehleranfällig, insbesondere beim Aktivieren innerhalb von Meetings. Unklar ist zudem, ob alle Teilnehmenden die App installiert haben müssen oder ob eine einmalige Freigabe ausreicht.

Inhaltlich erzeugt der Bot bereits während des Meetings kontinuierliche Zusammenfassungen. Dabei zeigt sich, dass die Qualität mit zunehmender Meetingdauer steigt: Anfangs dominiert Fließtext mit geringer Relevanz, später entwickelt sich eine strukturiertere Darstellung mit Stichpunkten, Überschriften und Hervorhebungen. Positiv ist, dass der Bot bestehende Inhalte nicht nur anhängt, sondern offenbar überarbeitet und konsolidiert. Negativ fällt auf, dass aktuell nahezu alles zusammengefasst wird – inklusive irrelevanter Kurzäußerungen („hm“, „ok“, Technikchecks). Die Relevanzfilterung fehlt noch.

Technisch wurde eine Message-Queue implementiert, um redundante oder sehr kurze Äußerungen nicht direkt an das LLM weiterzugeben. Ein offener Bug: Beim Beenden des Calls verbleiben unter Umständen Nachrichten in der Queue, die nicht mehr verarbeitet werden. Diese müssen beim Shutdown zwingend noch analysiert und in die Zusammenfassung integriert werden. Zusätzlich gab es unerklärliche Halluzinationen bzw. Fehlzuordnungen (z. B. Erwähnung von „Amazon“) sowie Zensur einzelner Begriffe, was vermutlich bereits auf Transkriptions- oder Microsoft-Ebene geschieht, nicht im LLM selbst.

Für die Weiterentwicklung wurden klare inhaltliche Leitplanken definiert. Der Prompt muss gezielt angepasst werden, um:
– irrelevante Inhalte konsequent auszublenden
– eine klare Meeting-Protokoll-Struktur zu erzwingen (Themen, Entscheidungen, To-Dos, offene Punkte)
– Relevanz über Vollständigkeit zu priorisieren

Die bisherige Einschätzung: Das System funktioniert technisch und konzeptionell, liefert aber aktuell noch keinen stabilen Mehrwert ohne Prompt-Tuning und Relevanzlogik. Mit klaren Strukturvorgaben hat der Ansatz jedoch eindeutig Potenzial.

Im Kontext der Masterarbeit wurde besprochen, dass:
– Meetings auf ca. 15 Minuten begrenzt werden, um Vergleichbarkeit zu sichern
– Echtzeit-Zusammenfassungen mit nachträglichen Zusammenfassungen verglichen werden (u. a. über BERTScore / ROUGE)
– Der Fokus der Dokumentation auf Architektur, Message-Queue, Filterlogik und Prompt-Design liegt, nicht auf umfangreichen Code-Snippets
– Screenshots des Frontends sinnvoll und erwünscht sind
– Code-Snippets nur dort eingesetzt werden sollen, wo sie konzeptionell etwas erklären (z. B. Filterlogik, Queue-Design)

Nächste konkrete Schritte:
– Prompt gezielt auf Relevanz und Struktur umbauen
– Queue-Handling beim Call-Ende fixen
– Installation/Startverhalten des Bots in Teams klären
– Weitere kurze Meetings (z. B. Daily) aufzeichnen und auswerten
– Quantitativen Vergleich Echtzeit- vs. Offline-Zusammenfassung durchführen

Unterm Strich: Die Basis steht, die Richtung ist richtig. Jetzt entscheidet sauberes Prompt-Engineering und klare Relevanzlogik darüber, ob aus einem technischen Prototyp ein echtes Produktmerkmal wird.

\section{Daily Stand Up 1 Meeting}
\subsection{Meetingtranskript}
\label{app:Daily1_transkript}
Daily Standup-20260210\_084627-Besprechungstranskript
10. Februar 2026, 07:46AM
17 Min. 0 Sek.

Yan Ohainski Transkription gestartet

Annika Mainitz   0:05
Sonja überall Henry mit drin.

Henri Reumschüssel   0:09
Sonja ist jetzt über Henry mit drin, ja.

Annika Mainitz   0:12
Dann wird er natürlich verwirrt sein.

Henri Reumschüssel   0:15
Oh, das wird ja stimmt. Das wird natürlich der Summary Bot nicht so gut finden. Oh, soll ich alleine rein?

Annika Mainitz   0:19
Mhm.

Yan Ohainski   0:28
Also, ein paar Sekunden habt ihr noch, bis der hier reinkommt. Das dauert immer ein bisschen.

Henri Reumschüssel   0:30
Sonst sag ich halt alles, ne?

Yan Ohainski   0:38
Gedacht, da ist er.

Annika Mainitz   0:41
Oh, how do some report?
Schön, dass du auch dabei bist. Wir freuen uns.

Henri Reumschüssel   0:46
Oh, ich freue mich auch so sehr, dass der Summary **** da ist.

Sonja Fiehn   0:49
Mhm.

Yan Ohainski   0:50
Schön, dass du auch dabei bist, A.

Sonja Fiehn   0:51
Äh.

Annika Mainitz   0:53
But.

Sonja Fiehn   0:53
Hallo.

Yan Ohainski   0:54
Mit mehr Leuten ist das ja richtig toll.

Joshua Sukowski   0:55
Morgen.

Yan Ohainski   0:58
Morgen.

Sonja Fiehn   1:14
Die Transkription wurde gestartet und Datenschutzrichtlinien sind schon verlinkt. Das ist ja wild.

Annika Mainitz   1:21
Oh.
Wild.

Sonja Fiehn   1:24
Ah, warte, zu Microsoft Datenschutz.

Yan Ohainski   1:30
Kommt Luca noch mit dazu?

Annika Mainitz   1:33
Das ist 'ne gute Frage. Keine Ahnung, der sitzt gerade im Auto auf dem Weg nach Wolfsburg.
Wir können ihn anfragen

Sonja Fiehn   1:59
Sonst kann ich schon mal nebenbei eh starten.
Umm.
Ich habe heute Vormittag ein Meeting zur AI Policy, dass das irgendwie so angeht. Und Henry, ich wär dir sehr dankbar, wenn wir uns wirklich zeitnah dann mit der DSFA auseinandersetzen können. und die fertig bekommen. Das haben wir gestern nicht geschafft.

Henri Reumschüssel   2:16
Ja.

Sonja Fiehn   2:22
Und dann würde ich gerne noch die.
Den ersten Schwung unternehmen für die Cold Calls. Ähm, fertig raussuchen. Da bin ich aber fast durch. Also ich glaube, es dauert noch so eine halbe Stunde oder so und dann bin ich glaube ich durch mit. Göttingen.
Und dann ist irgendwann die Frage, ob Luca da schon die.

Luca Mainitz   2:47
Luca.

Sonja Fiehn   2:47
Die Skripte vorbereitet hat.

Luca Mainitz   2:50
Nein.

Annika Mainitz   2:51
Oh.

Sonja Fiehn   2:52
I come.

Luca Mainitz   2:54
Aber die Skripte, die wir aktuell haben, sind auch gar nicht so schlecht. Also, da kann man auch einfach sich reinlesen.

Sonja Fiehn   2:59
OK.

Luca Mainitz   3:01
Um.

Sonja Fiehn   3:01
Wahrscheinlich würde ich mir das sowieso so ein bisschen personalisieren.

Luca Mainitz   3:06
Also ehrlich gesagt, ich glaube auch, nehmt euch mal das Base Script, lest euch die mal durch und dann probiert die mal aus und dann werdet ihr da irgendwas ***** finden und dann iteriert die halt.

Sonja Fiehn   3:17
Ja.

Luca Mainitz   3:18
Also, ich kannte über ja.

Sonja Fiehn   3:18
Genau, weil man muss ja irgendwie so seinen eigenen Touch reinbringen und dann immer sich ein bisschen ans Unternehmen anpassen und dann keine Ahnung.

Luca Mainitz   3:25
Ja, du, also ich habe doch die ersten paar Male das Spaceflip genutzt und dann irgendwann habe ich ich habe jetzt nur noch so eine kleine Tabelle, wo ich die Fragen habe, die ich abarbeiten will, während einem größeren Gespräch und dann ist man hier nur noch das.

Sonja Fiehn   3:37
Okay.

Luca Mainitz   3:37
Hello.
Wann kommt da der Rhythmus rein?

Sonja Fiehn   3:41
Ja.

Luca Mainitz   3:43
Yeah.

Sonja Fiehn   3:48
Ja.
Soweit zu meinen Plänen und dann haben wir heute Nachmittag ja wieder komplett Meeting Block.

Luca Mainitz   3:56
Mhm.
Apropos, heute Nachmittag Meeting, habt ihr gesehen, was ich gestern Nacht euch noch geschickt habe?

Sonja Fiehn   4:03
Ja, aber noch nicht angeschaut.

Annika Mainitz   4:04
No.

Luca Mainitz   4:04
Ich habe OK. Also, ich habe mich gestern Abend nochmal hingesetzt, weil.
Also, was Vittorio gestern gesagt hat, war ja nochmal sehr.
Der Film, wenn wir das Design finalisieren, ist es finalisiert. Dann haben wir es abgenickt, dann wird es auch nicht mehr richtig geändert. Es hat nicht aufhorchen lassen, weil momentan haben sie eigentlich unsere unsere Reviews da nicht so krass einfließen lassen.

Annika Mainitz   4:23
Mhm.

Luca Mainitz   4:30
Ja, schon ein bisschen, also vor allem das mit der Farbe hat sich jetzt deutlich verbessert.
deutlich schicker fand ich jetzt, aber da ist doch noch irgendwie viel improvement gewesen, funktionsmäßig. Und obwohl ich denen das, ich hatte denen extra nochmal eine E-Mail geschrieben, wo das alles drin stand, was ich jetzt letztes Mal auch noch kritisiert hab, was sie halt irgendwie ignoriert haben.
Und ich wollte jetzt noch mal unser Dokument lesen und das kann ich euch nur mal empfehlen und das Dokument lesen, was ich denen geschickt hatte, mit den ganzen korrigierten Anforderungen für User Stories. So hast du das gesehen, ne?

Sonja Fiehn   5:04
Mhm.

Luca Mainitz   5:05
dass
Das wollte ich alles nochmal durchlesen, weil ich glaube, die Anforderungen meiner Meinung sind eigentlich ganz klar definiert gewesen, aber die Designerin glaube ich, versteht die nicht. Und dann habe ich gedacht, was müssen wir machen? Wir müssen die Workflows aufzeichnen. Und dann habe ich gestern mal versucht, die generellen Workflows aufzuzeichnen, wie man was machen kann und was soll.
und das halt auch von der U.I., also so, wenn wir die U.I. angeguckt und gesagt, O.K., was muss ich wie machen und was, welche Elemente sind irgendwie wofür da und was soll man damit machen können und wieso, wieso können sie jetzt nicht und wie transition halt zwischen
Seitenleiste und Overlay. Wann switche ich? Wieso also wann switche ich, wenn ich switche? Was kriege ich für ein Mehrfaches oder nicht? Ja, dann hört es euch an, weil ich hab ihr habt den ganzen Tag keinen Zugriff auf mich und wir hören uns dann um 18:00 Uhr und dann will ich Input von euch.

Sonja Fiehn   5:44
K, nicht zu detailliert werden, aber ja, mhm.

Luca Mainitz   5:55
Deswegen und guckt euch das an und denkt es halt komplett einmal durch.
Das müsstet ihr auch nochmal machen. Ihr müsst jeden Schritt einmal euch durchdenken, weil wenn wir das nicht beim nächsten Design-Review fertig haben, dann wird es nicht gebaut und dann ist es so und dann kriegen wir es nicht wieder gefixt.

Sonja Fiehn   6:04
Ja.

Luca Mainitz   6:14
Und die wollten die Designreview diese Woche machen.

Sonja Fiehn   6:14
Weil sonst muss man da halt immer mit denen sprechen, weil ich bin halt auch echt nicht so happy mit dem Workflow allgemein mit Sierra, ne, weil ich finde es nicht so cool, wenn wir da halt vorher nichts geschickt bekommen. Ich kann mir das nicht anschauen, ich bin nicht in dem Meeting drin, ich habe da also und ich möchte auch was bei den Designs mitdiskutieren dürfen. Also deswegen finde ich

Luca Mainitz   6:29
Auf jeden Fall, auf jeden Fall.

Sonja Fiehn   6:31
Den Workflow insgesamt nicht OK. So, und dann können wir ja auch einfach mit denen drüber sprechen und denen das sagen, dass wir das halt anders möchten. Da würde ich mich ungerne komplett an.

Luca Mainitz   6:41
Ja, können wir auch noch. Können wir ja nochmal zusätzlich machen. Aber meine Meinung, also mein Plan wäre jetzt quasi gewesen, wir husteln das heute Abend noch mal komplett durch.

Sonja Fiehn   6:41
Da anpassen.

Luca Mainitz   6:50
Dann schreiben wir da irgendwann mal mega E-Mail da quasi zu, so nach dem Motto: "Hier, wir haben euch ja mal alle Workflows aufgezeichnet." Deswegen, wir werden es, ich würde, ich würde vorschlagen, wir diskutieren heute halt dieses Workflow-Chart noch mal durch. Ihr habt euch das vorher alles schon angeguckt, im besten Fall schon verändert, dass wir dann nur noch sagen können: "Hier, wir haben noch mal alle Workflows aufgemalt. Jetzt kann man es eigentlich nicht mehr falsch verstehen."
bitte stellt sicher, dass diese Wortfluss abgebildet sind. Ne, dann haben wir halt ein Dokument, was wir denen geben können, wo wir sagen können, das ist das, was wir auf jeden Fall brauchen und so soll es funktionieren. Und dann kann die Designerin sich überlegen, wie es aussehen soll dabei.

Sonja Fiehn   7:21
OK.
OK.
Das klingt doch eigentlich nach einem Plan.

Luca Mainitz   7:31
I.
Ich glaube nämlich auch, aber ja, ich bin mit dir, dass es nicht geschickt hat. Das nervt. Ich glaube, die Projektmanagerin ist *******.

Sonja Fiehn   7:39
Mhm.
Na gut, was steht noch bei euch an, Partypieps?

Annika Mainitz   7:51
Also, obviously das Design reviewen. Ich wollte noch weitere Beratungsunternehmen in Hannover raussuchen für unsere lustige Akquise.
Ähm.

Sonja Fiehn   8:06
Aber das läuft okay gerade für dich. Ja bis.

Annika Mainitz   8:08
Ja, ich habe 1000000 Tabs offen. Genau, und dann wurschtel ich mich da mal durch und gucke, wer jetzt tatsächlich passt und wer nicht. Und dann, du hast die im C. H. M. irgendwie hinterlegt oder wo hast du die gesammelt?

Sonja Fiehn   8:11
I can.
Yeah.
Mhm.

Annika Mainitz   8:24
OK.
OK, ja.

Luca Mainitz   8:26
Und Annika, dann guck dir nochmal an, wie du die hinterlegen musst. Also Sonja hat die jemals Kunde und Kalt hinterlegt.

Annika Mainitz   8:34
Mhm.

Luca Mainitz   8:34
Oder lauwarm.

Sonja Fiehn   8:34
Meistens, manchmal auch als lauwarm, und dann hab ich in Kommentar reingeschrieben, warum lauwarm.

Luca Mainitz   8:37
It's the line.
Ja, ich habe aber auch schon gesehen, dass du einen doppelt eingeflickt hast, den wir schon hatten, zum Beispiel.

Sonja Fiehn   8:45
Ah ja, super.

Luca Mainitz   8:47
kein Problem, hab deinen gelöscht, hab ihn neu angelegt und Kai, hab den existierenden genommen, weil der war besser als Eintrag. Ich konnte ihn nicht mergen, frag ich mich wieso und hab ihn dann als lauwarm markiert, weil wir ihn ja entscheiden schon kannten irgendwoher.

Sonja Fiehn   8:49
Meinen Sie das?
Ja.
Alright, ja, das mit dem Merch müssen wir uns noch mal anschauen, weil es gab auch sonst ein paar doppelte, glaube ich.

Luca Mainitz   9:06
Mhm.
Das funktioniert manchmal manchmal nicht und ich weiß nicht wieso.

Sonja Fiehn   9:12
Hm.

Luca Mainitz   9:12
Aber das funktioniert schon immer, manchmal, manchmal nicht.

Sonja Fiehn   9:15
Hm.
Na gut.

Annika Mainitz   9:22
Hmm.

Luca Mainitz   9:23
Oh, ich muss uns mal an.
Annika, kannst du das sonst einmal bitte auf mein Jira Board bauen, dass ich mir das angucken muss?

Annika Mainitz   9:40
Ui, dein Giraboard ist ganz schön voll.

Luca Mainitz   9:42
Mhm.
Thank you don't remind me just put it on there.

Annika Mainitz   9:47
OK, sorry.
Ist drin.

Luca Mainitz   9:53
Thank you.

Sonja Fiehn   10:15
Hallo und sonst so?

Henri Reumschüssel   10:19
Is ja ne.

Annika Mainitz   10:21
Mhm.

Henri Reumschüssel   10:25
Ich habe heute tatsächlich ausnahmsweise mal nicht so super viele super dringende Sachen zu tun.
Werden ein bisschen was aufarbeiten, was sonst noch so.
Herumliegt vielleicht ein bisschen was dokumentieren.
Genau.

Luca Mainitz   10:40
Das habe ich gerade nicht gehört. Entschuldigung, kannst du das bitte nochmal wiederholen?

Sonja Fiehn   10:44
Er hat gesagt, ein bisschen Win Olympische Winterspiele schauen.

Luca Mainitz   10:47
Das Ah, OK, Dann hab ich es auch richtig verstanden. OK.

Henri Reumschüssel   10:48
Ja, das wahrscheinlich auch, ja.

Luca Mainitz   10:54
Ne Henry, da wär natürlich Gangster, wenn du in das ****** mal nochmal reinguckst, ne?
Dass ich jetzt nicht irgendwas vergessen habe als Feature, was ihr schon gebaut habt.

Sonja Fiehn   10:59
Add you.

Luca Mainitz   11:03
und gerne reorganisiert das, schreibt da drum rum, arbeitet da wirklich einfach rein, weil irgendwie gestern, als wir das angefangen haben so durchzudiskutieren, ist mir danach klar geworden, ja, alles schön und gut, dass wir es Sierra gegeben haben, aber momentan haben wir es nicht niet und nagelfest dokumentiert und dann
Gerne, gerne nochmal irgendwie reingucken, dass es nie nagelfest ist.
Oder Sachen auch sagen brauchen wir nicht. Kannst du halt auch sagen, dann lieber fokussieren auf das, was wir brauchen.

Joshua Sukowski   11:41
O.K., ja, ich würd heute versuchen, irgendwie so weit wie möglich dieses Feature Flag Ding abzuschließen. Ich hab da gestern noch so 'n bisschen bisschen was von gemacht und funktioniert jetzt schon, dass man so
Teil Teile der Funktion dann halt ausblenden kann, wenn man die nicht sehen will oder den Leuten nicht zeigen möchte. Ja genau, das jetzt nochmal durchsprechen, gleich mit Henry, was genau da alles so rein muss und was so ein bisschen unnötig ist, was ich da schon gebaut habe oder so.
Genau.

Sonja Fiehn   12:18
K, das ist doch cool.

Yan Ohainski   12:30
Ja, ich bin noch weiter am Schreiben und werte gleich das Protokoll aus, was hier nebenbei gerade erstellt wird.

Sonja Fiehn   12:39
Mhm.

Yan Ohainski   12:39
Sieht auch ganz gut aus. Kann ich euch später im Römer zeigen oder allgemein einfach verschicken, damit ihr euch das mal anschauen könnt.

Annika Mainitz   12:47
Nice.

Sonja Fiehn   12:47
Mhm.

Yan Ohainski   12:49
Ein paar Sachen sind wirklich sehr wörtlich drin. Also, da stand einmal, dass Annika eine Millionen Tabs offen hat.

Joshua Sukowski   12:49
Hört sich gut an.

Annika Mainitz   12:56
Aha, perfekt. Ja, das ist ein wichtiger Punkt.

Luca Mainitz   12:58
Aber ja, gerne.
Schick das mal rum, wär ja spannend.

Yan Ohainski   13:06
Ja, seid ihr noch im Büro, Henry und Sonja?

Sonja Fiehn   13:10
Nee, noch nicht.

Yan Ohainski   13:12
Also, oder geht ihr noch ins Büro? So gefragt, ich bin ja auch noch nicht da.

Sonja Fiehn   13:17
Haben wir uns noch nicht abgestimmt? Wir haben kein Mittagessen. Also, ich denke mal.

Yan Ohainski   13:21
Ich hab auch nichts. Wir könnten zu Taco Fresh.

Sonja Fiehn   13:25
Mhm.

Annika Mainitz   13:25
Oh.

Luca Mainitz   13:26
Was? Das ist doch so geil. O.K., ich muss mich dann ausklinken, Leute. Wir hören uns heute Nachmittag. Bis dann, Tschüss.

Sonja Fiehn   13:33
Jo, alles klar.

Annika Mainitz   13:35
Mhm.

Sonja Fiehn   13:35
Bis dann, ciao ciao.

Joshua Sukowski   13:36
bis dann.

Annika Mainitz   13:37
Bye.

Sonja Fiehn   13:40
Henry, was denkst du?

Henri Reumschüssel   13:44
Umm.

Sonja Fiehn   13:44
Ach so Joshua, kannst du auch gleich natürlich schon ausklinken, wenn wir jetzt nur noch über Mittagessen sprechen.

Joshua Sukowski   13:50
Ich glaube, ich höre mir das an. Ich lasse mich mal ein bisschen inspirieren. Hier, ich habe auch noch nichts zum Mittagessen.

Sonja Fiehn   13:51
Auch wenn es natürlich. Ja, ne, verstehe.
Mhm.

Annika Mainitz   13:59
Tragt ihr mir dann meinen Taco Fresh Burrito hoch?

Sonja Fiehn   14:03
Oh.

Yan Ohainski   14:04
Hochtragen.

Annika Mainitz   14:07
Oder fahren? So, der meint auch Pakobrisch. Ich hab auch nichts zu essen.

Yan Ohainski   14:08
Mhm.

Sonja Fiehn   14:15
And home.

Yan Ohainski   14:20
Oh, hier steht Sonja, du hast das bestätigt.
Jetzt müssen wir den hochbringen.

Annika Mainitz   14:25
Ha, geil, danke somebody got. Ah, Ehrenechse.

Sonja Fiehn   14:26
Nein.
Ich möchte doch gerne die. Nein, nein, halt, stopp.
Oh oh.
Henry, magst du dich noch äußern, aber?

Henri Reumschüssel   14:37
können wir machen.

Sonja Fiehn   14:40
OK.

Yan Ohainski   14:41
Also dann ins Büro und Takko Fresh, okay?

Sonja Fiehn   14:46
All right.

Yan Ohainski   14:48
Ab wann wärt ihr denn da?
Jetzt gleich demnächst oder erst ein bisschen später? Oder wollen wir uns zum Mittag treffen oder wie?

Henri Reumschüssel   14:54
Na dann halt zu Mittag, würde ich sagen.

Sonja Fiehn   14:55
Ja, gerne dann zu Mittag, ehrlicherweise.

Yan Ohainski   14:58
Uns direkt dort treffen oder?
Erst im Büro.

Sonja Fiehn   15:04
Wir können uns eigentlich auch direkt dort treffen.

Yan Ohainski   15:07
Alright, dann 12:00 Uhr Takko Fresh.

Sonja Fiehn   15:09
12:00 Uhr Taco Fresh.

Henri Reumschüssel   15:11
Huh.

Yan Ohainski   15:12
Ai ai.

Sonja Fiehn   15:13
Mhm.
Bye bye.

Yan Ohainski   15:16
Okay, mach's gut. Bis dann.

Annika Mainitz   15:16
Viel Spaß mit euren Burritos. Bye bye.

Henri Reumschüssel   15:18
Ciao.

Sonja Fiehn   15:19
Dankeschön, ganz süß, mhm.

Joshua Sukowski   15:22
denn?

Yan Ohainski Transkription beendet

\subsection{Live-Zusammenfassung}
\label{app:Daily1_live}
**Teilnehmer:** Yan Ohainski, Annika Mainitz, Sonja Fiehn, Henri Reumschüssel, Joshua Sukowski, Luca Mainitz  

**Aktueller Stand \& Aufgaben:**  
- **Skripte \& Lead‑Management:** Luca arbeitet am Base‑Script, hat Duplicate‑Leads bereinigt und Lead‑Status (Kunde, Kalt, lauwarm) korrekt hinterlegt; Sonja hat das Vorgehen bestätigt.  
- **Cold Calls \& Recherche:** Sonja startet nach ca. 30 Minuten Recherche den ersten Cold‑Call‑Durchlauf. Luca plant ein gemeinsames „Durchhusteln“ am Abend und klärt den Rhythmus für weitere Calls.  
- **Design \& Workflow:** Erste UI‑Skizzen und Workflow‑Analyse (Seitenleiste, Overlay, Switch‑Timing, Mehrfach‑Anzeige) liegen vor. Team‑Review des aktuellen Workflow‑Diagramms ist für heute vorgesehen. Annika sucht ein Beratungs‑Unternehmen in Hannover für das Design‑Review.  
- **AI‑Policy \& DSFA:** Sonja informierte über das AI‑Policy‑Meeting und bat Henri, die Datenschutz‑Folgenabschätzung zeitnah abzuschließen (Henri zugestimmt).  
- **Film‑/Design‑Finalisierung:** Luca meldet verbesserte Farbgestaltung; funktionale Defizite bestehen noch. Das Film‑Konzept ist nach Design‑Finalisierung weitgehend fertig.  
- **Feature‑Flag (Joshua Sukowski):** Ziel ist, das Feature‑Flag‑Feature heute so weit wie möglich abzuschließen; weitere Abstimmung mit Henri zu notwendigen UI‑Elementen und „hide‑bar“ steht aus. Dokumentation und Funktionsdefinition stehen im Fokus.  
- **Koordination \& Termine:**  
  - Luca schlägt ein gemeinsames Gespräch um 18:00 Uhr vor.  
  - Henri ist heute besonders verfügbar und bestätigt seine Teilnahme.  
  - Lunch‑Absprachen: Henri schlägt ein Mittagessen vor, Sonja stimmt zu; Yan favorisiert ein Treffen im Büro, um Takko Fresh zu besorgen.  
  - Annika hat das neue Jira‑Ticket auf dem Board vermerkt.  
-**Protokoll \& Dokumentation (Yan Ohainski):** Yan arbeitet weiter am Protokoll, das später im Römer‑Meeting oder per Versand verteilt wird.  

**Kurzmeldungen \& aktuelle Entscheidungen:**  
- Annika bat um Unterstützung für ihren Taco Fresh Burrito („hochtragen“); Yan bestätigte, Sonja lehnte ab.  
- Joshua, Annika und Sonja gaben bestätigende Rückmeldungen („nice“, „hört sich gut an“, „mhm“).  
- Yan machte humorvoll darauf aufmerksam, dass Annika „eine Millionen Tabs offen hat“.  

**Verabschiedung (neue Informationen):**  
- Henri verabschiedete sich mit „Ciao.“  
- Sonja bedankte sich („Dankeschön, ganz süß, mhm.“).  
- Joshua stellte eine kurze Rückfrage („denn?“).  
- Annika wünschte allen „Viel Spaß mit euren Burritos“ und sagte „Bye bye.“

\subsection{LLM-Zusammenfassung}
\label{app:Daily1_llm}
Das Team ist vollständig gestartet, die automatische Transkription lief parallel und wird später noch ausgewertet und geteilt. Datenschutz-Hinweise wurden kurz thematisiert.

Schwerpunkt Organisation \& Rechtliches:
Ein zentrales Thema ist die AI Policy. Dazu gibt es am Vormittag ein separates Meeting. Parallel besteht hoher Druck, die DSFA zeitnah abzuschließen. Das ist offen geblieben und wird als klarer To-do-Punkt priorisiert.

Schwerpunkt Sales \& Akquise:
Die Vorbereitung der Cold Calls ist weit fortgeschritten. Die Zielunternehmen für Göttingen sind fast vollständig recherchiert (ca. 30 Minuten Restaufwand).
Es gibt kein neues finales Skript, aber Einigkeit, dass das bestehende Base-Script ausreichend ist. Vorgehen: lesen, ausprobieren, personalisieren, iterieren.
CRM-Pflege läuft, dabei wurden Dubletten entdeckt. Mergen funktioniert nicht zuverlässig und muss separat geklärt werden. Leads werden sauber als „kalt“, „lauwarm“ oder „Bestandskontakt“ markiert, inkl. Kommentaren.

Schwerpunkt Design \& Produkt (kritisch):
Das Design-Review ist aktuell der größte Risikopunkt. Klare Ansage:
– Wenn das Design einmal final abgenickt ist, wird es nicht mehr grundlegend geändert.
– Bisher fließen Reviews und Anforderungen nur unzureichend ein.
– Anforderungen sind dokumentiert, werden aber offenbar nicht richtig verstanden.

Konsequenz:
Es werden konkrete, durchgängige Workflows (User-Flows) erstellt – unabhängig vom visuellen Design. Fokus liegt auf Funktion, Übergängen, Zuständen und Logik (Sidebar vs. Overlay, Wechsel, Mehrfachauswahl etc.).
Diese Workflows werden heute vorbereitet, intern durchdacht und heute Abend gemeinsam reviewt. Ziel ist ein eindeutiges Dokument, das nicht mehr fehlinterpretiert werden kann. Das Design soll sich danach richten, nicht umgekehrt.
Zusätzlich besteht Unzufriedenheit mit dem aktuellen Abstimmungsprozess mit dem externen Design-Team (keine Vorab-Unterlagen, eingeschränkte Mitwirkung). Das soll offen adressiert werden.

Schwerpunkt Tech:
– Feature Flags: Umsetzung ist weit fortgeschritten. Funktionen können gezielt ein- und ausgeblendet werden. Heute Abstimmung, was wirklich gebraucht wird und was raus kann.
– Dokumentation: Bestehende Features und Anforderungen sollen nochmal kritisch geprüft, ergänzt oder bewusst gestrichen werden, um Fokus zu halten.

Weitere Aufgaben:
– Weitere Beratungsunternehmen in Hannover für Akquise recherchieren und korrekt im CRM hinterlegen.
– Ein Jira-Ticket zur CRM-Problematik (Dubletten/Merge) wurde angelegt.
– Transkript/Meetingprotokoll wird später geteilt.

Kurz gesagt:
Operativ läuft vieles gut, aber Design und Anforderungen müssen jetzt hart nachgeschärft werden. Wenn die Workflows nicht diese Woche sauber stehen, werden falsche Dinge gebaut – und das wäre teuer. Heute Abend ist dafür ein entscheidender Punkt.

\section{Daily Stand Up 2 Meeting}
\subsection{Meetingtranskript}
\label{app:Daily2_transkript}
Daily Standup-20260211\_084941-Besprechungstranskript
11. Februar 2026, 07:49AM
13 Min. 12 Sek.

Yan Ohainski Transkription gestartet

Luca Mainitz   0:04
*****.
Das bedeutet, ich krieg Kuchen und Kaffee.

Sonja Fiehn   0:09
ha ha ha ha.

Luca Mainitz   0:09
Hab ich wieder glücklich.
Dicke 3 Stunden, Junge, brutal oder?

Annika Mainitz   0:12
Das hört sich nett an.

Sonja Fiehn   0:16
Klingt nach einem feinen Frühstück.

Luca Mainitz   0:19
Aber ich komme so ein bisschen später, also jetzt.
Mhm.

Annika Mainitz   0:23
Hat konntest du wenigstens schlafen, hat sich Moritz zusammengerissen.

Luca Mainitz   0:28
Nein, also ja, er war früh im Bett, aber der ******* hat halt nicht gepennt, weil er Podcasts angehabt und sich die ganze Zeit hin und her geworfen hat.

Annika Mainitz   0:35
aus

Luca Mainitz   0:35
Und dann gab es oho packs um 10 nach 12 und dann konnte ich auch direkt schlafen.

Henri Reumschüssel   0:41
Mhm.

Sonja Fiehn   0:43
Mhm.

Annika Mainitz   0:43
Oh Mann, ja, das ist ein bisschen nervig.

Luca Mainitz   0:46
Unsere Sleeprice Energy sind wirklich Menschen nicht.
Also, ich mag schlafen nicht.

Sonja Fiehn   0:53
Mhm.

Luca Mainitz   0:55
Very market schlafe. So it was.

Sonja Fiehn   1:00
Mhm.

Luca Mainitz   1:00
Ich habe gestern noch den Branch fertig gemacht für das WhatsApp Integration Ding, aber ich habe das noch nicht end-zu-ende testen können.
Das ist das, was ich heute machen wollte.

Sonja Fiehn   1:14
Nice.

Luca Mainitz   1:15
Einfach aus Lust und Laune.

Annika Mainitz   1:23
Hm.

Sonja Fiehn   1:24
Hm.
Ja.
Ich würde mich mal wieder mit Pitchdecks beschäftigen, also Live Pitchdeck für HTI.
Und written Pitcheck für Andreas.
Ähm.
Wird bei dem H.W.G.K. Forschungstag Viech anrufen.
Um.
Und.
Ich weiß, dann dann mal schauen, was noch so Prio 1 Themen sind.
Ja.

Luca Mainitz   2:02
Aber wir kriegen heute mit, ob wir dritten oder draußen sind.

Sonja Fiehn   2:04
Oh.
Wie bitte?

Annika Mainitz   2:06
Mhm.

Luca Mainitz   2:08
Ich sing halt mit, ob wir drinnen oder draußen sind.

Annika Mainitz   2:09
ATI.

Sonja Fiehn   2:11
Ja.

Luca Mainitz   2:14
Da muss ich mal ganz leid, dass Sie drin sind.

Annika Mainitz   2:15
Yep.

Sonja Fiehn   2:17
Hast du mal.

Annika Mainitz   2:18
Ja, ich hoffe doch. Wäre asozial, sonst, wenn er halt einfach geschrieben hätte, so nach dem Motto "Ja, cooler Pitch" und dann so "Ja, aber nicht cool genug Leute übrigens".

Sonja Fiehn   2:26
Mhm.

Luca Mainitz   2:27
Den Rick? Nein, der hat uns geschrieben, weil wir glaub ich schon da rausgejagt haben. Das hat er erst nicht gemacht.

Annika Mainitz   2:34
Ist eigentlich so.
Die Hoffnung und weil er ja weiß, dass wir auch eine längere Anfahrt hätten, dass man das vielleicht auch planen muss, ne?

Luca Mainitz   2:36
Um.
Hast du ja auch schon fast gedacht, dass das so ein bisschen hin war. So, hey, piepst euch mal vor.
Aber egal, heute wissen wir mehr, was ist denn unser Battleplan? Dann müssen wir ja eigentlich ab morgen Nachmittag das erste Mal Pitch üben, oder?

Sonja Fiehn   2:56
Mhm.

Annika Mainitz   2:58
Na dann ist den Rest der Woche eigentlich Abriss, Patch, Verarbeitung. Alles andere kann warten.

Luca Mainitz   2:58
Dann danach eigentlich.
Ja, und dann jeden Tag, ja.

Sonja Fiehn   3:05
Wollte mir ehrlich gesagt ganz gerne Freitag freinehmen und dafür dann am Wochenende mehr arbeiten.

Luca Mainitz   3:05
Meinst du?
Das kannst du auch machen.

Sonja Fiehn   3:11
Aber ich kann hab dann trotzdem am Freitag auch zwischendurch Zeit, meinen Kram durchzuüben.

Annika Mainitz   3:19
Mhm.

Luca Mainitz   3:22
Also, ich würde Samstag, Samstag, Sonntag müssten wir eigentlich nur irgendwie von dir einmal was üben oder irgendwie uns noch einmal treffen, vielleicht Sonntag. Beziehungsweise, die Frage ist, fahren wir denn Montag alle zusammen mit einem Auto?

Sonja Fiehn   3:30
Mhm.

Luca Mainitz   3:39
Hat seine ne.

Henri Reumschüssel   3:39
Ja, das sind viele, ne.

Luca Mainitz   3:40
Ja, denk ich auch. Und dann können wir mit dem Auto schon die ganze Zeit immer Flitschen üben. Das klingt ein bisschen dumm, aber
Genau, hast du drei Stunden Zeit, die man hat, so? Oh my not.

Sonja Fiehn   3:50
Mhm.
Two.
Aber das wäre der Pitch, wäre am Dienstag, ne?

Luca Mainitz   3:59
Ach, Dienstag erst, ich dachte Montag, Annika, du hast erst Montag gesagt.

Sonja Fiehn   4:01
Oder, oder Montag.

Annika Mainitz   4:02
Nee, Montag ist das 16.

Sonja Fiehn   4:04
Oh, talk.

Luca Mainitz   4:07
Guck nochmal bitte jetzt auf der Website nach, please.

Annika Mainitz   4:07
Also bin ich mir ziemlich sicher.

Sonja Fiehn   4:10
Dann müssen wir halt ja am 16. am Montag. Das heißt im Worst Case, wenn es früh ist, müsste man natürlich auch schon Sonntag fahren.

Luca Mainitz   4:10
Please triple check.
Just.
Und hat Annika auch gesagt, meinte sie so: 'Boah, eigentlich die 3 Stunden kannst du auch immer morgens abreißen.'
Also, vor 9 kann ich es mir nicht vorstellen.

Annika Mainitz   4:34
Ja, aber wenn es um neun wäre, dann müssten wir ja schon um halb sechs losfahren.
Damit wir Puffer haben.

Luca Mainitz   4:41
Jetzt sind Sie eher Marke Ungeil, aber jetzt nichts irgendwie komplett außerhalb dieser Welt.
oder

Annika Mainitz   4:48
Also, ich weiß nicht, ob wir dann so.

Luca Mainitz   4:48
Also keine anderen.
Ah.
Gegen alle Männer schläft man so gut in einem Hotel.
Dann können wir also dann diskutieren.

Sonja Fiehn   5:03
Hm.

Annika Mainitz   5:06
Habt ihr gerade die E-Mail gesehen, die Mandy uns geschrieben hat?

Sonja Fiehn   5:11
Mhm.

Annika Mainitz   5:11
Das verwirrt mich komplett.

Henri Reumschüssel   5:11
Mhm.

Annika Mainitz   5:14
Wo kommen diese 64.000 her?

Luca Mainitz   5:19
Die wir noch haben.

Annika Mainitz   5:23
Ja, anscheinend, aber wo sollen die herkommen? Das verwirrt mich vollends.

Luca Mainitz   5:29
Anrufen.
Wir haben die förderliche Richtlinie mal wieder zu den Grünen geändert.

Annika Mainitz   5:48
Komplett verwirrt. Wo sollte dieses Geld herkommen?

Luca Mainitz   5:53
OK, peeps.

Annika Mainitz   5:53
Hä?

Sonja Fiehn   5:53
Vielleicht auch noch der Hälter von uns, mhm.

Luca Mainitz   5:56
Let's no mal kurz here back to the roots.
Umm.
Weil ich muss jetzt gleich rein. Was, was steht an bei euch? Was steht nicht an bei euch?

Annika Mainitz   6:11
Anscheinend werde ich Mandy anrufen.
Mhm.

Henri Reumschüssel   6:20
Ich werd erst mal reindiven und schauen, ob ich irgendwie reproduzieren kann, was Marvin da für ein Problem hat.
Das sieht aber extrem weird aus.
Na ja.

Luca Mainitz   6:36
Hey, ist es ein Technikleiter auf unserer Seite?

Henri Reumschüssel   6:40
Nicht direkt unsere Seite, aber wir müssen ja damit umgehen.

Luca Mainitz   6:45
Ja, dann drin ist es kein.

Henri Reumschüssel   6:47
Genau, also wir müssen eine Lösung dafür finden, weil es wohl sowohl in der Office-Web-App als auch in der Windows-App, also in den meistbenutzten Apps, vorkommt.

Luca Mainitz   6:56
Ja, ist vorher aber nicht vorgekommen, oder also?
Das irgendwie irritiert, fand ich.

Henri Reumschüssel   7:00
Ja, kann sein. Ging aber jetzt schon länger bei ihm, meinte er. Er konnte das immer wieder reproduzieren.

Luca Mainitz   7:06
OK.

Sonja Fiehn   7:14
OK.
Jan, du schreibst wahrscheinlich.
Weiter einfach.

Yan Ohainski   7:24
Genau.

Sonja Fiehn   7:26
Ich würde eigentlich gerne gucken, dass ich dir, boah, ich hoffe, ich kann dir das heute noch schicken, die mit denen deine bis einschließlich Kapitelmethoden. Das wäre wahrscheinlich ganz gut, ne?

Yan Ohainski   7:32
Mhm.
Ja, ich wär ich bin auch schon soweit und hab die nächsten Kapitel ready.

Sonja Fiehn   7:43
Mhm.

Henri Reumschüssel   7:44
Ich hab bis Implementierung auch schon durchgelesen, also bis Ende Implementierung.

Yan Ohainski   7:47
Ja, das hat sie am Wochenende. Das habe ich auch schon eingearbeitet. Oder hattest du das noch weitergelesen?

Henri Reumschüssel   7:52
Ich hab gestern weitergelesen.

Yan Ohainski   7:53
Ach so, ja, OK. Dann schaue ich da auch nochmal rein und arbeite das ein. Danke schön.

Henri Reumschüssel   7:58
Und es Ergebnisse kann man das auch schon lesen, oder?

Yan Ohainski   8:01
Ergebnisse, ja.
Also, ich hab das jetzt schon mal vorgeschrieben. Ich muss jetzt noch das von heute einarbeiten, aber ich nehme mal an, dass sich das nicht großartig unterscheidet.

Henri Reumschüssel   8:04
OK.

Yan Ohainski   8:11
Von den Grundsätzen Ergebnissen.

Henri Reumschüssel   8:14
Okay.

Yan Ohainski   8:15
Heute mach ich dann noch Diskussion und schlussfertig und dann steht soweit der gesamte Text und dann nur noch überarbeiten.

Sonja Fiehn   8:16
Ich bin jetzt mit.
OK.
Okay.

Henri Reumschüssel   8:26
So, und was macht der Summary Bot, das faule Schwein?

Yan Ohainski   8:29
Der Summary Bot ist ordentlich am Summarizing.

Sonja Fiehn   8:33
Und was heißt das?

Yan Ohainski   8:34
Aber qualitativ nicht so hochwertig, aber er fasst etwas zusammen.
Ich hatte nur am Anfang vergessen, die Transkription zu starten, und dann ist die Qualität immer sehr schlecht. Der versucht das immer irgendwie ins Englische zu bringen.
Yo.

Annika Mainitz   8:50
Hm.

Yan Ohainski   8:53
Aber Auswertung kommt gleich.

Annika Mainitz   8:57
Nice.
Ich warte auch auf die Zusammenfassung von gestern. Das würde mich ja auch brennend interessieren. Sehr gut.

Yan Ohainski   9:02
Oh ja, stimmt. Ja, die kann ich dir schicken.
Ja.

Luca Mainitz   9:05
Endlich.

Henri Reumschüssel   9:05
Ach so, die wichtigen Leute haben schon bekommen.

Yan Ohainski   9:08
die Leute, die bei mir saßen, im Beruf.

Luca Mainitz   9:08
Langsame Frage.

Annika Mainitz   9:08
*****.

Henri Reumschüssel   9:10
Ihre
Aber die ist witzig. Annika war auf jeden Fall. Hat Jan eigentlich bestätigt, dass er dir einen Burrito hochbringt? Hat der Summary Bot behauptet?

Annika Mainitz   9:21
Ja, OK, du schuldest mir einen Burrito.

Yan Ohainski   9:24
Ich schuld den Burrito in der in der Zusammenfassung steht Sonja hat das bestätigt, glaub ich.

Luca Mainitz   9:24
Oh my god.

Annika Mainitz   9:26
Yeah.
Hm.

Henri Reumschüssel   9:29
Nein, ich dachte, das stimmt.

Sonja Fiehn   9:29
Nein, ich glaube aber nicht.

Yan Ohainski   9:30
Oh, okay, dann habe ich das falsch verstanden.

Sonja Fiehn   9:31
Das wurde wieder, wurde wieder verbessert.

Yan Ohainski   9:35
Gar nicht. Der kann sich nicht selbst verbessern. Der wird nur schlechter mit der Zeit.

Annika Mainitz   9:41
Hm.

Sonja Fiehn   9:42
Aus deiner Perspektive vielleicht. Da vielleicht mal einen unabhängigen Gutachter rüberjagen.

Yan Ohainski   9:48
Aber es ist auch teilweise serviert, was für Transkripte Teams hier erstellt. Also, da sind auch sehr viele Fehler drin. Sehr viele sind auch einige Fehler drin, die natürlich 1 zu 1 übernommen werden.

Henri Reumschüssel   9:57
Ja, die Qualität ist schon mies, ne, von den Teams-Transkripten.

Luca Mainitz   10:00
Ja, das ist glaub ich das erste, was man fixen müsste.

Yan Ohainski   10:02
No.
Aber es ist ja auch aktuell nur die Notlösung. Das ist ja eigentlich ein Plan, das anders zu machen mit eigenen Datenquellen und nicht mit dem, was Microsoft ja bei sich auf den Servern dann auch hat.

Luca Mainitz   10:09
Ja.
Ah, OK, good.

Yan Ohainski   10:17
Aber das grundsätzliche Prinzip.

Henri Reumschüssel   10:22
Alles klar, dann haben wir es.

Sonja Fiehn   10:27
Dit war dit hö.

Annika Mainitz   10:27
Ja, Luca wollte noch eine dumme Frage stellen, glaube ich.

Luca Mainitz   10:31
Bei Jan, ne?

Henri Reumschüssel   10:32
Das glaube ich gerne.

Yan Ohainski   10:34
Was?

Henri Reumschüssel   10:35
Was?

Luca Mainitz   10:36
Oh, I had me got sassy.

Sonja Fiehn   10:37
Was?

Henri Reumschüssel   10:40
heute

Luca Mainitz   10:42
Umm.
Ne, ich wollte einfach, ich wollte fragen, kannst du nicht quasi so die ersten 5 Sachen, die so transkribiert werden, durchjagen und sagen, erkenne die Sprache? Und dann stellt er sich danach selbst auf die Sprache, weil du damit den Systemraum modifizierst.

Yan Ohainski   10:57
Nee, das ist ja eine Teams-eigene Sache. Wenn du keine Transkription von Teams startest, dann nimmt der Standardsachen.

Luca Mainitz   11:01
Ach, die Teams Transkriptionssprache.
Ah, OKOKOK ja.

Yan Ohainski   11:06
Und ich hab schon geguckt, du kannst nicht über das Frontend, in dem ich das jetzt gemacht habe, auf diese Funktion zugreifen, um sie selbst manuell zu starten, nicht manuell, um sie mit dem Bot zu starten.

Luca Mainitz   11:15
Kriegt man denn jetzt den? Kriegt man jetzt den Road Audio Stream? Kriegt man den?

Yan Ohainski   11:19
Den kannst du an sich auch mit dir holen.
Ja.

Luca Mainitz   11:22
OK.

Yan Ohainski   11:24
Hab ich aber jetzt aktuell, hab ich einmal ausprobiert.

Luca Mainitz   11:24
Ja, dann ist das nächste da einfach irgendwie sich nen bezahlten Transcription Service zu holen aus Deutschland den allein zu jagen, damit das besser wird, dann hast du ja schon erstmal 50% gefixt.

Yan Ohainski   11:33
Ja, das geht. Das wäre auf jeden Fall eine Option.

Luca Mainitz   11:36
Also, so alles klar.

Yan Ohainski   11:40
Alrighty, ich komme dann zum Mittagessen ins Büro mit Essen.

Annika Mainitz   11:44
Ganz gut.

Henri Reumschüssel   11:46
Super, bis dann.

Yan Ohainski   11:47
OK, dann bis später.

Sonja Fiehn   11:47
Bist auch so ein Mitesser. Bis dann, genau. Ciao, ciao.

Luca Mainitz   11:47
But like, ciao.

Yan Ohainski   11:50
Ciao ciao.

Annika Mainitz   11:50
Tschüss, bis dann.

Henri Reumschüssel   11:52
Oh.

Summary Bot Transkription beendet


\subsection{Live-Zusammenfassung}
\label{app:Daily2_live}
- **Sonja Fiehn**  
  - Gibt kurze Bestätigungen („Mhm“, „OK“, „Weiter einfach“, „Okay“).  
  - Will Freitag freinehmen, dafür am Wochenende länger arbeiten; hat trotzdem etwas Zeit am Freitag zum Üben.  
  - Möchte heute noch Unterlagen zu den Kapitel‑Methoden (bis einschließlich) senden.  
  - Signalisiert, dass sie „mit“ ist – also auf dem aktuellen Stand (bestätigt erneut).  
  - Äußert Unsicherheit: „Nein, ich glaube aber nicht.“ (zu nicht näher genannter Frage).  
  - Regt an, einen unabhängigen Gutachter aus Sicht des Gegenübers hinzuziehen.  

- **Luca Mainitz**  
  - Lockerer, teils füllender Sprachstil.  
  - Klärt Missverständnis zu Rick, verweist auf die erneut geänderte Förder‑Richtlinie der Grünen.  
  - Plant Pitch‑Training ab morgen Nachmittag; danach täglicher Übungs‑Ablauf.  
  - Schlägt zusätzliches Wochenende‑Training (Samstag oder vorzugsweise Sonntag) vor.  
  - Während der Autofahrt zum Pitch am **Montag, 16.** sollen kontinuierlich „Flitschen“ (Impro‑Übungen) stattfinden.  
  - Kann **nicht vor 9 Uhr** trainieren; bei Start um 9 Uhr ist Abfahrt **5:30 Uhr** (Puffer nötig).  
  - Neue Aufgabe: Website prüfen und alle Informationen **dreifach** verifizieren.  
  - Vorschlag, nachzurufen, um offene Fragen zu klären.  
  - Fragt im aktuellen Batch nach dem Aufgaben‑Stand: **„Was steht an bei euch? Was steht nicht an bei euch?“**  
  - Kurz bestätigende Rückmeldungen („OK“, „Okay“, „Mhm.“).  
  - Hat bis zum Ende der Implementierung alles durchgelesen; Ergebnisse sind bereits einsehbar.  
  - Bestätigt erneut seine aktuelle Bereitschaft; macht im Batch humorvollen Kommentar zum Summary‑Bot.  
  - Hinweis, dass die wichtigen Personen die Unterlagen bereits erhalten haben („Ach so, die wichtigen Leute haben schon bekommen.“).  
  - Humorvolle Anmerkungen: „Aber die ist witzig. Annika war auf jeden Fall.“  
  - Fragt, ob Jan bestätigt hat, einen Burrito zu bringen, und ob der Summary Bot etwas behauptet hat.  
  - **Neue Information:** Äußert Vertrauen mit dem Satz „Das glaube ich gerne.“  
  - **Neuer Hinweis aus aktuellem Batch:** Verabschiedet sich mit „Super, bis dann.“  

- **Annika Mainitz**  
  - Kurze Zustimmungen, zeigt Verwirrung („Hä?“, „Komplett verwirrt…“).  
  - Bestätigt Pitch‑Tag **Montag, 16.** und dass das dreistündige Training **morgens** stattfinden kann.  
  - Hinweis auf **ATI** und Warnung vor laienhaftem Pitch‑Feedback sowie fehlender Einbindung passender Personen.  
  - Plant, Mandy anzurufen (z.B. wegen der fehlenden E‑Mail).  
  - Fokus der restlichen Woche: **Abriss, Patch und Verarbeitung**; andere Themen können warten.  
  - Fragen zur Herkunft einer Geldsumme (z.B. **64000**) und nach einer **E‑Mail von Mandy**.  
  - Wartet auf die gestrige Zusammenfassung („Ich warte auch auf die Zusammenfassung von gestern…“) und möchte diese dringend erhalten.  
  **Neuer Hinweis aus aktuellem Batch:** Gibt ein kurzes, zustimmendes „Ganz gut.“  

- **Henri Reumschüssel**  
  - Kurz bestätigende Rückmeldungen („OK“, „Okay“, „Mhm.“).  
  - Hat bis zum Ende der Implementierung alles durchgelesen; Ergebnisse sind bereits einsehbar.  
  - Bestätigt erneut seine aktuelle Bereitschaft; macht im Batch humorvollen Kommentar zum Summary‑Bot.  
  - Hinweis, dass die wichtigen Personen die Unterlagen bereits erhalten haben („Ach so, die wichtigen Leute haben schon bekommen.“).  
  - Humorvolle Anmerkungen: „Aber die ist witzig. Annika war auf jeden Fall.“  
  - Fragt, ob Jan bestätigt hat, einen Burrito zu bringen, und ob der Summary Bot etwas behauptet hat.  
  - **Neue Information:** Äußert Vertrauen mit dem Satz „Das glaube ich gerne.“  
  - **Neuer Hinweis aus aktuellem Batch:** Verabschiedet sich mit „Super, bis dann.“  


- **Yan Ohainski**  
  - Knapp bestätigende Rückmeldungen („Genau“, „Mhm.“).  
  - Hat bereits Vorentwürfe erstellt, muss heute noch die aktuellen Inhalte einarbeiten.  
  - Plant, heute die Diskussion abzuschließen, Text schlussfertig zu machen; anschließend bleibt nur noch Überarbeitung.  
  - Erwähnt Ergebnisse zu den Grundsätzen (Prinzipien‑Ergebnisse) und bestätigt, dass diese vorliegen.  
  - Zu Beginn des Tages vergaß er, die Transkription zu starten, was zu schlechter Qualität führte.  
  - Das System versucht häufig, Inhalte ins Englische zu übersetzen, liefert jedoch nur mittelmäßige Zusammenfassungen.  
  - Erkennt nun Fehler im Verständnis: „Oh, okay, dann habe ich das falsch verstanden.“  
  - Bemängelt, dass viele im Team erstellte Transkripte zahlreiche Fehler enthalten, die häufig 1 zu 1 übernommen werden.  
  - **Neue Informationen:**  
    - Sagt, dass er zum Mittagessen ins Büro kommt und Essen mitbringt („Alrighty, ich komme dann zum Mittagessen ins Büro mit Essen.“).  
    - Bestätigt eine Option mit „Ja, das geht. Das wäre auf jeden Fall eine Option.“.  
  - **Zusätzliche Aussage:** Betont, dass die aktuelle Lösung nur eine Notlösung ist und künftig eigene Datenquellen statt Microsoft‑Server‑Daten genutzt werden sollen.  
  - **Weitere neue Details:** Stellt mehrfach die Frage „Was?“ (Verwirrung).

\subsection{LLM-Zusammenfassung}
\label{app:Daily2_llm}
Das Team startet mit lockerem Austausch, geht dann aber zügig in operative Themen.

Luca hat den Branch zur WhatsApp-Integration fertiggestellt. End-to-End-Tests stehen noch aus und sind für heute geplant. Das ist aktuell sein Hauptfokus.

Sonja arbeitet an zwei Pitchdecks: einem Live-Pitchdeck für HTI sowie einer schriftlichen Version für Andreas. Zusätzlich plant sie einen Anruf wegen des Forschungstags und priorisiert danach weitere Themen. Freitag möchte sie frei nehmen und stattdessen am Wochenende arbeiten, bleibt aber teilweise erreichbar.

Zum HTI-Pitch steht heute die finale Rückmeldung an, ob das Team dabei ist. Falls ja, wird ab morgen konsequent Pitch-Training gestartet. Plan ist tägliches Üben. Logistisch wird diskutiert, ob gemeinsame Anreise oder ggf. Hotelübernachtung sinnvoll ist, abhängig vom Termin (Montag oder Dienstag, wird noch final geprüft).

Annika klärt offene Fragen zu einer E-Mail von Mandy bezüglich unerwarteter 64.000 € Fördermitteln – Herkunft unklar, daher direkter Klärungsanruf.

Henri analysiert ein reproduzierbares Problem aus der Office-Web- und Windows-App, das bei Marvin auftritt. Ziel ist, Ursache einzugrenzen und Lösung zu finden. Technisches Risiko besteht, da es mehrere Kernplattformen betrifft.

Yan arbeitet intensiv an seiner schriftlichen Ausarbeitung: Ergebnisse stehen, heute folgen Diskussion und Schluss, danach nur noch Überarbeitung. Parallel entwickelt er den Summary-Bot weiter. Aktuell leidet die Qualität unter fehlerhaften Teams-Transkripten. Mittelfristig ist geplant, von der Microsoft-Transkription auf bessere bzw. eigene Audio-/Transkriptionsquellen oder externe Services umzusteigen. Technisch möglich wäre auch der direkte Zugriff auf den Raw-Audio-Stream.

Kurz zusammengefasst die Prioritäten:

Heute: WhatsApp-Tests, Pitchdecks, Fördermittelklärung, Bug-Reproduktion, Dokument fertigstellen.
Diese Woche: Pitch-Vorbereitung intensivieren, sobald Teilnahme bestätigt.
Risiken: unklare Finanzierung, instabile Transkriptqualität, App-Bug.